\chapter{Linear Algebra}\label{ch:linear}

\begin{center}\itshape
Why is $F_{n+1}F_{n-1}-F_n^2$ always $\pm1$? Why is the determinant of $\left(\frac1{i+j-1}\right)_{i,j}$ positive, though tiny?\\
Matrices answer questions that do not mention matrices.
\end{center}

Cut an $8\times8$ chessboard into four pieces and put the pieces back as a $5\times13$ rectangle. The square has $64$ cells. The rectangle has $65$. Where is the extra cell?

The extra cell is a gap. Along the diagonal of the rectangle the pieces leave a long thin parallelogram of area exactly $1$ (Figure~\ref{fig:chessboard}). Its longer side has length $\sqrt{73}$, so the altitude on that side is $1/\sqrt{73}$, a little under an eighth of a cell. A small drawing hides it. The numbers $5$, $8$ and $13$ are consecutive Fibonacci numbers, and
\[
5\cdot13-8^2=1.
\]
That is one case of the identity in the epigraph. Jean-Dominique Cassini, then director of the Paris Observatory, stated it in 1680. Robert Simson published a proof in 1753. A version of the dissected board was found among the papers of Lewis Carroll. The identity is the determinant of a $2\times2$ matrix, and the gap is that determinant read as an area.

Vector spaces are where that picture lives, and linear maps are what move the pieces. Rank and nullity count what a map forgets and what it keeps. The determinant is the signed volume that makes Cassini's $\pm1$ an area. A change of basis then asks when a map is just a stretch along each axis. An inner product chooses those axes to be perpendicular: that is the spectral theorem. The adjugate turns the same determinant into a formula for the inverse. Least squares is that same perpendicular, used on three points that miss a single line.

Determinants were found twice, about ten years apart. In 1683 Seki Takakazu wrote a manuscript on eliminating unknowns; the $3\times3$ determinant is there, with a rule for the signs. In 1693 Leibniz described the same quantity in a letter to l'H\^opital, and his notation for the entries is already close to our $a_{ij}$. The letter was not printed until 1850. Cramer gave his rule for linear systems in 1750. Cauchy, in a memoir of 1812, defined the general determinant by the signs of permutations, proved that the determinant of a product is the product of the determinants, and fixed the name. Sylvester supplied the word \emph{matrix} in 1850. Cayley, in a memoir of 1858, treated matrices as things one adds and multiplies. He proved that a $2\times2$ matrix satisfies its own characteristic equation and stated the general case, adding that he would not write out the proof for every size. The proof for every $n$ is Frobenius's, from 1878, and it is Theorem~\ref{thm:cayleyhamilton} below. Hamilton had met a quaternion form of the same idea in 1853. Cauchy's theorem that a real symmetric form has real eigenvalues is from 1829, and it was a theorem about the axes of an ellipsoid: the matrix language was not there yet. The Hilbert matrix of the epigraph is from an 1894 paper on approximation by polynomials. An eigenvector can also rank pages on the web. Brin and Page computed one in 1998, as students at Stanford. Bryan and Leise later titled an article ``the twenty-five-billion-dollar eigenvector''. That title is theirs. A ranking can be a direction fixed by a matrix.

\begin{figure}[ht]
\centering
\fig{ch07-chessboard}
\caption{The $64=65$ paradox. Left: an $8\times8$ square cut into two trapezia and two triangles along lines through lattice points. Right: the four pieces rearranged into a $5\times13$ rectangle. The ``diagonal'' of the rectangle is not a line: the edges of the pieces have slopes $\frac25$ and $\frac38$, and between them lies a thin parallelogram (shaded, drawn to scale) with vertices $(0,0)$, $(8,3)$, $(13,5)$, $(5,2)$. Its area is $\left|\det\begin{pmatrix}8&5\\3&2\end{pmatrix}\right|=|16-15|=1$: the missing square, smeared along the diagonal.}
\label{fig:chessboard}
\end{figure}

\section{Vector spaces, from scratch}

An arrow in the plane can be slid, stretched, and added to another arrow by the parallelogram rule. The same two operations, addition and scaling, make sense for triples of numbers, for polynomials, and for sequences. A vector space is the list of rules those operations obey. We write them once, and then every later argument may quote them.

\begin{definitionx}[vector space]
Let $\mathbb K$ be $\R$ or $\C$. A \emph{vector space} over $\mathbb K$ is a set $E$ with vector addition and scalar multiplication. Addition is associative and commutative, there is a zero vector $0$, and every vector has an additive inverse. For vectors $v,w$ and scalars $a,b$, the scalar rules are
\[
a(v+w)=av+aw,\quad (a+b)v=av+bv,\quad (ab)v=a(bv),\quad 1v=v.
\]
These rules imply $0v=0$ and $a0=0$: from $0v=(0+0)v=0v+0v$ subtract $0v$, and from $a0=a(0+0)=a0+a0$ subtract $a0$.
\end{definitionx}

In $\mathbb K^n$ both operations act one coordinate at a time. Polynomials, functions on a fixed set, and sequences are vector spaces with the same pointwise operations. A \emph{subspace} is a nonempty subset closed under addition and scalar multiplication. It inherits the rules. The plane $x+y+z=0$ inside $\R^3$ is a subspace, because both operations preserve the equation. The plane $x+y+z=1$ is not: it does not contain $0$.

A \emph{linear combination} of $v_1,\dots,v_p$ is $\sum_i a_iv_i$. Their \emph{span} is the set of all such combinations. They are \emph{linearly independent} when $\sum_i a_iv_i=0$ forces every $a_i=0$. They \emph{span} $E$ when every vector of $E$ is a combination. A \emph{basis} has both properties, so every vector has exactly one expression in it. The coefficients are its \emph{coordinates}. A space is \emph{finite-dimensional} when some finite family spans it. The empty family is a basis of $\{0\}$, and that space has dimension $0$.

On the plane $x+y+z=0$, the vectors $(1,0,-1)$ and $(0,1,-1)$ are a basis. Every vector there is $x(1,0,-1)+y(0,1,-1)$, and the first two coordinates force the coefficients. For polynomials of degree at most $d$, the monomials $1,X,\dots,X^d$ are a basis, because two polynomials are equal only when their coefficients match.

Figure~\ref{fig:bases} is the same fact in the plane. The arrow $v=(3,1)$ does not change. In the standard basis its coordinates are $3$ and $1$. In the basis $f_1=(1,1)$, $f_2=(1,0)$ they are $1$ and $2$, since $f_1+2f_2=(3,1)$. A basis is a choice of rulers along the arrow.

\begin{figure}[ht]
\centering
\fig{ch07-bases}
\caption{One vector, two bases. The charcoal arrow is $v=(3,1)$ in both panels. Left: $v=3e_1+e_2$ on the rust and ochre axes. Right: the same arrow is $f_1+2f_2$, with $f_1=(1,1)$ and $f_2=(1,0)$. The dashed sides are that sum. Coordinates belong to the basis; $v$ does not.}
\label{fig:bases}
\end{figure}

The word ``dimension'' is safe only because of the next theorem. It says that every basis of a finite-dimensional space has the same number of vectors. Everything later leans on that count.

\begin{theoremx}[exchange lemma, and dimension]
Let $E$ be spanned by $v_1,\dots,v_n$.
\begin{enumerate}[label=(\alph*)]
\item \textbf{Steinitz's exchange lemma}: if $u_1,\dots,u_p$ are linearly independent in $E$, then $p\le n$.
\item $E$ has a basis, extracted from $v_1,\dots,v_n$; every linearly independent family extends to a basis; and all bases have the same number of elements, called $\dim E$.
\item A family of $\dim E$ vectors that is linearly independent, or that spans $E$, is a basis. Every subspace $F$ of $E$ is finite-dimensional with $\dim F\le\dim E$, and $\dim F=\dim E$ only if $F=E$.
\end{enumerate}
\end{theoremx}
\begin{proof}
(a) The independent vectors $u_i$ can replace the spanning vectors $v_j$ one at a time. Write $u_1=\sum_j\lambda_jv_j$. Some $\lambda_j\ne0$, since $u_1\ne0$. After renumbering, $\lambda_1\ne0$, and then $v_1$ is a combination of $u_1,v_2,\dots,v_n$, so this new family still spans $E$. Suppose that after $k<p$ steps the family $u_1,\dots,u_k,v_{k+1},\dots,v_n$ spans $E$. This already forces $k\le n$. Write
\[
u_{k+1}=\sum_{i\le k}\mu_iu_i+\sum_{j>k}\lambda_jv_j.
\]
If every $\lambda_j$ were $0$, the vectors $u_i$ would be dependent. So some $\lambda_j\ne0$, which forces $k<n$, and that $v_j$ can be exchanged for $u_{k+1}$. After $p$ steps we have exchanged $p$ of the $v_j$, so $p\le n$.

(b) Among the subfamilies of $v_1,\dots,v_n$ that span $E$, take one of smallest size. It is linearly independent: a dependence would let us drop one vector and still span. So it is a basis. If $u_1,\dots,u_p$ are independent, add, one by one, those $v_j$ that are not combinations of the vectors already chosen. The family stays independent and ends by spanning $E$. Finally, if $B$ and $B'$ are bases with $|B|=b$ and $|B'|=b'$, then (a) applied both ways gives $b\le b'$ and $b'\le b$.

(c) An independent family of $\dim E$ vectors extends to a basis, and a basis has $\dim E$ elements, so nothing was added. A spanning family of $\dim E$ vectors contains a basis, so nothing was removed. For a subspace $F$, independent families in $F$ have at most $\dim E$ elements by (a). A maximal one exists, and it is a basis of $F$: a vector of $F$ outside its span could be added. If $\dim F=\dim E$, that basis is a basis of $E$ by the first claim of (c), so $F=E$.
\end{proof}

\subsection*{Linear maps, kernel and image}

A map $u:E\to F$ between vector spaces over the same field is \emph{linear} when
\[
u(av+bw)=au(v)+bu(w)
\]
for all vectors and scalars. In particular $u(0)=0$, since $u(0)=u(0+0)=u(0)+u(0)$. Define
\[
\ker u=\{v\in E:u(v)=0\},\qquad \operatorname{Im}u=\{u(v):v\in E\}.
\]
Both are subspaces. The image of a combination is the combination of the images, and the same sentence with the image equal to $0$ shows that the kernel is a subspace. The map $u$ is injective exactly when $\ker u=\{0\}$: the equality $u(v)=u(w)$ is the equality $u(v-w)=0$. If $E$ is finite-dimensional, the images of a finite spanning family span $\operatorname{Im}u$. The dimension of the image is the \emph{rank} of $u$, written $\operatorname{rk}u$.

Differentiation on polynomials is the example to keep. Let $\R_d[X]$ be the polynomials of degree at most $d$, and let $D:\R_2[X]\to\R_1[X]$ be $D(P)=P'$. Then $D$ is linear, its kernel is the constant polynomials, and it is onto: $D\bigl(aX+\frac12 bX^2\bigr)=a+bX$. The kernel has dimension $1$ and the image has dimension $2$. Those add to $3$, the dimension of the domain. Figure~\ref{fig:rank} draws the count without coordinates. The kernel is what is collapsed. The image is what remains.

\begin{figure}[ht]
\centering
\fig{ch07-rank}
\caption{Rank--nullity. A linear map collapses a subspace, the kernel, and lands on the image. The two dimensions add to the dimension of the domain. Differentiation on polynomials of degree at most $2$ collapses the constants and hits every polynomial of degree at most $1$.}
\label{fig:rank}
\end{figure}

\begin{proposition}[rank--nullity]\label{prop:ranknullity}
For $u:E\to F$ linear with $\dim E<\infty$, $\dim\ker u+\operatorname{rk}u=\dim E$.
\end{proposition}
\begin{proof}
Take a basis $(k_1,\dots,k_p)$ of $\ker u$ and extend it to a basis $(k_1,\dots,k_p,e_1,\dots,e_r)$ of $E$, which the exchange theorem allows. The vectors $u(e_i)$ span $\operatorname{Im}u$, because the $u(k_i)$ are zero. They are independent: if $\sum\lambda_iu(e_i)=0$, then $\sum\lambda_ie_i$ lies in $\ker u$, so it is a combination of the $k_i$, and every coefficient is zero by independence of the basis. Thus $\operatorname{rk}u=r=\dim E-p$.
\end{proof}

If $E$ and $F$ have the same finite dimension, rank--nullity says that $u$ is injective exactly when it is surjective. A bijective linear map has a linear inverse. To see the inverse is linear, apply $u$ to both sides of the proposed rule $u^{-1}(ay+bz)=au^{-1}(y)+bu^{-1}(z)$. Both sides have the same image, and $u$ is injective, so the two sides agree.

\subsection*{Matrices record a linear map}

Choose a basis $(e_j)$ of $E$ and a basis $(f_i)$ of $F$. Put the coordinates of $u(e_j)$ into column $j$ of a matrix $A=(a_{ij})$. Linearity then says $[u(v)]=A[v]$: if $v=\sum_jx_je_j$, then
\[
u(v)=\sum_{i,j}a_{ij}x_jf_i.
\]
An $m\times n$ matrix is therefore a map $\mathbb K^n\to\mathbb K^m$. We write $M_{m,n}(\mathbb K)$ for these matrices, and $M_n(\mathbb K)$ when the matrix is square.

If two maps can be composed, their matrices multiply. On a coordinate vector, $A(Bx)=(AB)x$, where
\[
(AB)_{ik}=\sum_ja_{ij}b_{jk}.
\]
Associativity of the product is associativity of composition. Commutativity can fail. Take
\[
A=\begin{pmatrix}0&1\\0&0\end{pmatrix},\qquad B=\begin{pmatrix}0&0\\1&0\end{pmatrix}.
\]
Then $AB=\operatorname{diag}(1,0)$ and $BA=\operatorname{diag}(0,1)$. Turn a book onto its face and then onto its side, or do the two turns in the other order: the corners do not finish in the same place. These two matrices are that fact.

The identity matrix is the identity map. A square matrix is \emph{invertible} exactly when its map is bijective, and the inverse matrix is the inverse map. Its columns are independent exactly when its kernel is $\{0\}$, so rank--nullity is another test for invertibility.

A change of basis is also a matrix. If the columns of $P$ are the coordinates of a new basis $(e'_j)$ in the old basis, then $[v]_{\mathrm{old}}=P[v]_{\mathrm{new}}$. The coordinate maps are bijections, so $P$ is invertible. For a linear map $u:E\to E$ with old matrix $A$, the relation $AP[v]_{\mathrm{new}}=P[u(v)]_{\mathrm{new}}$ says that the new matrix is $P^{-1}AP$. Such matrices are called \emph{similar}. They describe one map in two languages.

Here is a small check. Let $u(P)(X)=P(2X)$ on $\R_1[X]$. In the basis $(1,X)$ the matrix is $A=\operatorname{diag}(1,2)$. In the basis $(1,1+X)$,
\[
P=\begin{pmatrix}1&1\\0&1\end{pmatrix},\qquad P^{-1}AP=\begin{pmatrix}1&-1\\0&2\end{pmatrix}.
\]
You can see the columns without multiplying: $u(1)=1$, and $u(1+X)=1+2X=-1+2(1+X)$. Similarity changes the description and leaves the map fixed.

\begin{examplex}[the $2\times2$ case, worked]
In the plane, the matrix $R_\theta=\begin{pmatrix}\cos\theta&-\sin\theta\\ \sin\theta&\cos\theta\end{pmatrix}$ sends $e_1=(1,0)$ to $(\cos\theta,\sin\theta)$ and $e_2$ to $(-\sin\theta,\cos\theta)$. It is rotation by $\theta$, and $R_\theta R_\varphi=R_{\theta+\varphi}$ by the addition formulas. Exercise~\ref{ex:twobytwo} reads the same identity as multiplication of complex numbers. The matrix $\begin{pmatrix}1&1\\0&1\end{pmatrix}$ fixes $e_1$ and sends $e_2$ to $e_1+e_2$. It is a \emph{shear}: each horizontal line slides along itself by its own height. The area does not change. Figure~\ref{fig:parallelogram} is that slide, and Figure~\ref{fig:orientation} is the sign. The matrix $\begin{pmatrix}a&b\\c&d\end{pmatrix}$ is invertible if and only if $ad-bc\ne0$, with inverse
\[
\frac1{ad-bc}\begin{pmatrix}d&-b\\-c&a\end{pmatrix}.
\]
Multiply the two matrices and see. The kernel is $\{0\}$ if and only if the columns are not proportional, which is the same condition $ad-bc\ne0$. Before the determinant has a name, it already decides invertibility, and it already measures the parallelogram.
\end{examplex}

\begin{figure}[ht]
\centering
\fig{ch07-orientation}
\caption{Signed area in the plane. Left: the rust edge $v$ then the ochre edge $w$, taken counterclockwise, give a positive determinant. Right: the same parallelogram with the edges swapped. The arc now runs clockwise, and the determinant changes sign. The absolute value is the area either way.}
\label{fig:orientation}
\end{figure}

\begin{figure}[ht]
\centering
\fig{ch07-parallelogram}
\caption{Left: the parallelogram spanned by the columns $v_1=(a,c)$ and $v_2=(b,d)$ of a matrix; its area is $|ad-bc|$. Middle: a shear, adding a multiple of $v_1$ to $v_2$, slides the top side along itself and changes neither the area nor the determinant (Theorem~\ref{thm:det}(d)). Right: after shearing $v_2$ into $w_2\perp v_1$, the parallelogram is a rectangle of area $\|v_1\|\cdot\|w_2\|$, and $\|w_2\|$ is the distance from $v_2$ to the line $\R v_1$.}
\label{fig:parallelogram}
\end{figure}

\portrait{sylvester}{James Joseph Sylvester (1814--1897). Photograph, nineteenth century. Public domain, via Wikimedia Commons.}

\section{Determinants}

The number $ad-bc$ has two jobs in the plane: it vanishes exactly on the singular matrices, and its absolute value is the area of the column parallelogram. The same number exists for every size. It is the unique function of the columns that is linear in each column, vanishes when two columns agree, and equals $1$ on the identity.

Let $\mathfrak S_n$ be the group of permutations of $\{1,\dots,n\}$.

\begin{lemma}[sign of a permutation]\label{lem:sign}
$\varepsilon(\sigma)=\prod_{i<j}\frac{\sigma(j)-\sigma(i)}{j-i}$ is $\pm1$ and satisfies $\varepsilon(\sigma\tau)=\varepsilon(\sigma)\varepsilon(\tau)$. Every transposition has sign $-1$, and every permutation is a product of transpositions.
\end{lemma}
\begin{proof}
The quotient $\frac{\sigma(j)-\sigma(i)}{j-i}$ does not change when $i$ and $j$ are swapped, so it depends only on the pair $\{i,j\}$, and $\varepsilon(\sigma)=\prod_{\text{pairs}}\frac{\sigma(j)-\sigma(i)}{j-i}$. As $\{i,j\}$ runs over the pairs, so does $\{\sigma(i),\sigma(j)\}$: numerator and denominator are the same product up to sign, so $\varepsilon(\sigma)=\pm1$. Next,
\[
\begin{aligned}\varepsilon(\sigma\tau)&=\prod_{\text{pairs}}\frac{\sigma(\tau(j))-\sigma(\tau(i))}{\tau(j)-\tau(i)}\\&\quad\cdot\prod_{\text{pairs}}\frac{\tau(j)-\tau(i)}{j-i}=\varepsilon(\sigma)\varepsilon(\tau),\end{aligned}
\]
the first product being $\varepsilon(\sigma)$ written over the pairs $\{\tau(i),\tau(j)\}$. For the transposition $(1\,2)$, only the pair $\{1,2\}$ gives a negative factor, so $\varepsilon=-1$. Any transposition $(a\,b)$ equals $\pi(1\,2)\pi^{-1}$ for some $\pi$, so it has sign $\varepsilon(\pi)(-1)\varepsilon(\pi)^{-1}=-1$. Finally, by induction on $n$: if $\sigma(n)=k\ne n$, then $(k\,n)\sigma$ fixes $n$ and is a product of transpositions of $\{1,\dots,n-1\}$.
\end{proof}

Concretely, $\varepsilon(\sigma)=(-1)^N$, where $N$ is the number of \emph{inversions}: pairs $i<j$ with $\sigma(i)>\sigma(j)$. Each inversion contributes one negative factor. The permutation $(3,1,2)$ of $\{1,2,3\}$, meaning $1\mapsto3$, $2\mapsto1$, $3\mapsto2$, has two inversions, so it is even. It is a $3$-cycle, a product of two transpositions. Swapping two columns of a matrix is a transposition. The sign is the reason that swap changes the determinant.

\begin{theorem}[determinant]\label{thm:det}
\begin{enumerate}[label=(\alph*)]
\item There is a unique function $\det:M_n(\mathbb K)\to\mathbb K$ that is linear in each column, vanishes when two columns are equal (\emph{alternating}), and satisfies $\det I_n=1$; it is $\det A=\sum_{\sigma\in\mathfrak S_n}\varepsilon(\sigma)\prod_{j=1}^na_{\sigma(j),j}$. Every multilinear alternating $D$ satisfies $D(A)=D(I_n)\det A$.
\item $\det(AB)=\det A\det B$ and $\det A^{\mathsf T}=\det A$; in particular everything said about columns holds for rows.
\item $A$ is invertible if and only if $\det A\ne0$.
\item Adding to a column a multiple of another does not change the determinant, and the determinant of a triangular matrix is the product of its diagonal.
\end{enumerate}
\end{theorem}
\begin{proof}
(a) Let $D$ be multilinear and alternating. Expanding $D(x+y,x+y)=0$ in two columns shows that $D$ changes sign when two columns are swapped. Write column $j$ of $A$ as $\sum_ia_{ij}e_i$ and expand by multilinearity:
\[
D(A)=\sum_{f}\prod_ja_{f(j),j}\,D(e_{f(1)},\dots,e_{f(n)}),
\]
the sum running over all maps $f:\{1,\dots,n\}\to\{1,\dots,n\}$. Terms with $f$ not injective vanish, because two columns agree. For $f=\sigma$ a permutation, write $\sigma$ as a product of $t$ transpositions and undo them one at a time: $D(e_{\sigma(1)},\dots,e_{\sigma(n)})=(-1)^tD(I_n)=\varepsilon(\sigma)D(I_n)$. This proves uniqueness, and the formula $D(A)=D(I_n)\det A$.

Conversely, the formula is multilinear: each term contains exactly one entry from each column. If columns $k\ne l$ are equal, let $\tau=(k\,l)$. The map $\sigma\mapsto\sigma\tau$ pairs the permutations. The products $\prod_ja_{\sigma\tau(j),j}$ and $\prod_ja_{\sigma(j),j}$ agree, because columns $k$ and $l$ agree, while $\varepsilon(\sigma\tau)=-\varepsilon(\sigma)$. The terms cancel in pairs. And $\det I_n=1$, because only $\sigma=\mathrm{id}$ contributes.

(b) For fixed $A$, the map $X\mapsto\det(AX)$ is multilinear and alternating in the columns of $X$, since column $j$ of $AX$ is $A$ times column $j$ of $X$. By (a) it equals $\det(A)\det X$. For the transpose, substitute $\sigma\mapsto\sigma^{-1}$ in the formula. Use $\varepsilon(\sigma^{-1})=\varepsilon(\sigma)$ and $\prod_ja_{\sigma(j),j}=\prod_ia_{i,\sigma^{-1}(i)}$.

(c) If $AA^{-1}=I_n$, then $\det A\det A^{-1}=1$. If $A$ is not invertible, its columns are dependent: $C_k=\sum_{j\ne k}\lambda_jC_j$ for some $k$. By linearity in column $k$, $\det A$ is a sum of determinants with two equal columns, hence $0$.

(d) The first claim is linearity in that column, together with alternation. For an upper triangular matrix, a nonzero term of the formula needs $\sigma(j)\le j$ for every $j$, which forces $\sigma=\mathrm{id}$ by induction from $j=1$. A lower triangular matrix is the transpose of an upper one.
\end{proof}

The sum has $n!$ terms. For $n=10$ that is $3{,}628{,}800$ terms, which is why nobody expands a large determinant that way. What one uses is (d): row operations bring the matrix to triangular form in about $n^3$ arithmetic operations. Laplace expansion, next, reduces one $n\times n$ determinant to $n$ determinants of size $n-1$. It is the right tool when a row is almost full of zeros. It is the wrong tool when nothing is zero.

The chessboard is now a one-line case of (b). Let $F_0=0$, $F_1=1$, and $F_n=F_{n-1}+F_{n-2}$ for $n\ge2$. Set $M=\begin{pmatrix}1&1\\1&0\end{pmatrix}$, so $\det M=-1$.

\begin{examplex}[Cassini, from a determinant]
For every $n\ge1$,
\[
M^n=\begin{pmatrix}F_{n+1}&F_n\\F_n&F_{n-1}\end{pmatrix}.
\]
The case $n=1$ is $M$ itself, since $F_2=1$ and $F_0=0$. If the formula holds for $n$, then
\[
M^{n+1}=M^nM=\begin{pmatrix}F_{n+1}+F_n&F_{n+1}\\F_n+F_{n-1}&F_n\end{pmatrix},
\]
which is the same shape at $n+1$. Theorem~\ref{thm:det}(b) gives $\det(M^n)=(\det M)^n=(-1)^n$. The left side is $F_{n+1}F_{n-1}-F_n^2$. For $n=6$, this is $13\cdot5-8^2=1$, the area of the gap in Figure~\ref{fig:chessboard}.
\end{examplex}

Exercise~\ref{ex:fibmat} asks for this identity again, and then for Binet's formula. Binet needs the eigenvalues of $M$, which are the next tool. The area is already settled.

\begin{propositionx}[Laplace expansion]
Let $A_{ij}$ be the $(n-1)\times(n-1)$ matrix obtained from $A$ by deleting row $i$ and column $j$, and $c_{ij}=(-1)^{i+j}\det A_{ij}$ the \emph{cofactor}. Then for every row $i$ and every column $j$,
\[
\det A=\sum_{k=1}^na_{ik}c_{ik}=\sum_{k=1}^na_{kj}c_{kj}.
\]
\end{propositionx}
\begin{proof}
Expand along column $j$. By linearity in that column, $\det A=\sum_ka_{kj}\det A^{(k)}$, where $A^{(k)}$ is $A$ with column $j$ replaced by $e_k$. Move column $j$ to the first position by $j-1$ swaps of adjacent columns, then move row $k$ to the first position by $k-1$ swaps of adjacent rows. The sign changes by $(-1)^{j-1+k-1}=(-1)^{k+j}$, and the matrix now has first column $e_1$ and, below and to the right of the corner, the matrix $A_{kj}$. In the formula of Theorem~\ref{thm:det}(a), the only nonzero terms have $\sigma(1)=1$, and they form the determinant of $A_{kj}$. So $\det A^{(k)}=(-1)^{k+j}\det A_{kj}$. The row expansion is the same argument applied to $A^{\mathsf T}$.
\end{proof}

\begin{examplex}
Along the first row,
\[
\begin{aligned}
\det\begin{pmatrix}2&1&3\\0&4&1\\1&2&5\end{pmatrix}
&=2\det\begin{pmatrix}4&1\\2&5\end{pmatrix}-\det\begin{pmatrix}0&1\\1&5\end{pmatrix}
+3\det\begin{pmatrix}0&4\\1&2\end{pmatrix}\\
&=2\cdot18-(-1)+3\cdot(-4)=25.
\end{aligned}
\]
Along the first column the same determinant is $2\cdot18+1\cdot(1-12)=25$. The signs $(-1)^{i+j}$ form a chessboard, with $+$ in the corner. Expand along the row or column with the most zeros.
\end{examplex}

\begin{think}
Compute $\det\begin{pmatrix}1&a\\1&b\end{pmatrix}$, then guess $\det\begin{pmatrix}1&a&a^2\\1&b&b^2\\1&c&c^2\end{pmatrix}$ without expanding. It is a polynomial in $a,b,c$ of total degree $3$, and it vanishes when two of the variables are equal.
\end{think}

\begin{proposition}[Vandermonde]\label{prop:vdm}
For $x_1,\dots,x_n\in\mathbb K$,
\[
\begin{aligned}V(x_1,\dots,x_n)&=\det\left(x_i^{\,j-1}\right)_{1\le i,j\le n}\eqbr=\prod_{1\le i<j\le n}(x_j-x_i).\end{aligned}
\]
\end{proposition}
\begin{proof}
Proceed by induction on $n$. The case $n=1$ is $V(x_1)=1$. If two of the $x_i$ are equal, both sides vanish. Otherwise, replace $x_n$ by a variable $X$ and expand along the last row. Then $P(X)=V(x_1,\dots,x_{n-1},X)$ is a polynomial of degree at most $n-1$, and the coefficient of $X^{n-1}$ is the cofactor of the corner entry, which is $V(x_1,\dots,x_{n-1})$. The polynomial $P$ vanishes at $x_1,\dots,x_{n-1}$, because two rows agree. By Proposition~\ref{prop:rootsfactor},
\[
P=V(x_1,\dots,x_{n-1})\prod_{i<n}(X-x_i).
\]
Evaluate at $X=x_n$ and apply the induction hypothesis.
\end{proof}

\portrait{vandermonde}{Alexandre-Th\'eophile Vandermonde (1735--1796). Portrait, eighteenth century. Public domain, via Wikimedia Commons.}

The Vandermonde determinant is why Lagrange interpolation (Theorem~\ref{thm:lagrange}) is a statement about an invertible matrix. The system ``find a polynomial of degree less than $n$ with prescribed values at $n$ distinct points'' has matrix $(x_i^{j-1})$. That matrix is invertible exactly when the points are distinct. Exercise~\ref{ex:circulant}, in the chapter on complex numbers, computes a circulant determinant by the same diagonalisation.

\section{Reduction}

A matrix is a linear map written in a basis. A good basis can make the matrix almost transparent: diagonal, when such a basis exists. The numbers on the diagonal are the \emph{eigenvalues}. The directions of the basis are the \emph{eigenvectors}. The map then stretches by $\lambda_i$ along the $i$th direction and does nothing else. The job of this section is to decide when that basis exists, and to find it.

Let $A\in M_n(\mathbb K)$ and $\chi_A(X)=\det(XI_n-A)$.

\begin{reminder}[eigenvalues, eigenvectors, similarity]
A scalar $\lambda\in\mathbb K$ is an \emph{eigenvalue} of $A$ if $Av=\lambda v$ for some $v\ne0$, an \emph{eigenvector}. The \emph{eigenspace} $\ker(A-\lambda I_n)$ collects them, together with $0$. Two matrices $A$ and $B=P^{-1}AP$, with $P$ invertible, are \emph{similar}: they are the same linear map written in two bases, and the columns of $P$ are the new basis vectors. The matrix $A$ is \emph{diagonalisable} if it is similar to a diagonal matrix, that is, if $\mathbb K^n$ has a basis of eigenvectors of $A$. A \emph{polynomial in} $A$ is $P(A)=\sum p_kA^k$ for $P=\sum p_kX^k$, with $A^0=I_n$. The map $P\mapsto P(A)$ respects sums and products, because the powers of $A$ commute with each other.
\end{reminder}

\begin{proposition}[characteristic polynomial]\label{prop:charpoly}
With $\operatorname{tr}A=\sum_i a_{ii}$, the \emph{trace}, $\chi_A=X^n-\operatorname{tr}(A)X^{n-1}+\dots+(-1)^n\det A$, and the eigenvalues of $A$ are its roots in $\mathbb K$. Over $\C$, $\operatorname{tr}A$ and $\det A$ are the sum and the product of the eigenvalues, counted with multiplicity as roots of $\chi_A$. Similar matrices have the same $\chi$.
\end{proposition}
\begin{proof}
In the formula for $\det(XI_n-A)$, a permutation $\sigma\ne\mathrm{id}$ moves at least two indices, so its term contains at most $n-2$ diagonal entries $X-a_{jj}$ and has degree at most $n-2$. The coefficients of $X^n$ and $X^{n-1}$ therefore come from $\prod_j(X-a_{jj})$: they are $1$ and $-\operatorname{tr}A$. The constant term is $\chi_A(0)=\det(-A)=(-1)^n\det A$. Next, $\lambda$ is an eigenvalue if and only if $\ker(\lambda I_n-A)\ne\{0\}$, which is $\det(\lambda I_n-A)=0$ by Theorem~\ref{thm:det}(c). Over $\C$, $\chi_A$ splits by Theorem~\ref{thm:dalembert}, and Vi\`ete's formulas (Proposition~\ref{prop:viete}) give the sum and the product of the roots. Finally
\[
\det(XI-P^{-1}AP)=\det\bigl(P^{-1}(XI-A)P\bigr)=\chi_A.
\]
\end{proof}

\begin{examplex}[a $2\times2$ matrix, diagonalised by hand]
For $A=\begin{pmatrix}2&1\\1&2\end{pmatrix}$, one has $\chi_A=(X-2)^2-1=(X-1)(X-3)$. For $\lambda=1$, the matrix $A-I=\begin{pmatrix}1&1\\1&1\end{pmatrix}$ kills $u_1=(1,-1)$. For $\lambda=3$, the matrix $A-3I=\begin{pmatrix}-1&1\\1&-1\end{pmatrix}$ kills $u_2=(1,1)$. Let $P$ have these eigenvectors as columns, $P=\begin{pmatrix}1&1\\-1&1\end{pmatrix}$. Then $AP=P\operatorname{diag}(1,3)$, so $A=PDP^{-1}$ and
\[
A^n=PD^nP^{-1}=\frac12\begin{pmatrix}3^n+1&3^n-1\\3^n-1&3^n+1\end{pmatrix}.
\]
Check the case $n=1$. The trace $4=1+3$ and the determinant $3=1\cdot3$ agree with Proposition~\ref{prop:charpoly}.

The matrix $N=\begin{pmatrix}0&1\\0&0\end{pmatrix}$ is the contrast. Its characteristic polynomial is $X^2$, so the only eigenvalue is $0$, and $\ker N=\mathbb Ke_1$ is only a line. There is no basis of eigenvectors, so $N$ is not diagonalisable. An eigenvalue is present. There are too few eigenvectors to span the plane.
\end{examplex}

\begin{think}
A projection satisfies $A^2=A$, and a reflection satisfies $A^2=I$. Both are diagonalisable. Is every matrix with $A^2=0$ diagonalisable? What distinguishes the polynomials $X^2-X$ and $X^2-1$ from $X^2$?
\end{think}

\begin{theorem}[diagonalisability]\label{thm:diag}
$A$ is diagonalisable if and only if $P(A)=0$ for some polynomial $P=\prod_{i=1}^r(X-\lambda_i)$ with distinct $\lambda_i\in\mathbb K$. In that case $\mathbb K^n=\bigoplus_i\ker(A-\lambda_iI_n)$: every vector is a unique sum of vectors from these eigenspaces. This is the meaning of a \emph{direct sum}.
\end{theorem}
\begin{proof}
If $A=QDQ^{-1}$ with $D$ diagonal, let $\lambda_1,\dots,\lambda_r$ be the distinct diagonal entries. Then $\prod_i(D-\lambda_iI)=0$, because each diagonal entry is killed by one factor, and therefore
\[
\prod_i(A-\lambda_iI)=Q\Bigl(\prod_i(D-\lambda_iI)\Bigr)Q^{-1}=0.
\]
Conversely, let $L_i$ be the Lagrange polynomials on the points $\lambda_1,\dots,\lambda_r$ (Theorem~\ref{thm:lagrange}). Interpolating the constant $1$ gives $\sum_iL_i=1$, so every $v\in\mathbb K^n$ is $v=\sum_iL_i(A)v$. The product $(X-\lambda_i)L_i$ is a constant multiple of $P$, so $(A-\lambda_iI)L_i(A)v=0$. Thus $L_i(A)v$ lies in $\ker(A-\lambda_iI)$, and the eigenspaces span $\mathbb K^n$. The sum is direct. If $\sum_iv_i=0$ with $Av_i=\lambda_iv_i$, apply $L_j(A)$. On $v_i$ it acts as multiplication by $L_j(\lambda_i)=\delta_{ij}$, so $v_j=0$. A basis adapted to the decomposition diagonalises $A$.
\end{proof}

The box above asked about $A^2=0$. The polynomial $X^2$ has a double root, so the theorem does not apply, and $N$ is a counterexample. The relations $A^2=A$ and $A^2=I$ give $X(X-1)$ and $(X-1)(X+1)$, split with simple roots. Projections are therefore diagonalisable over any field. Reflections are diagonalisable over any field in which $1\ne-1$.

Two consequences are used constantly. A matrix whose characteristic polynomial has $n$ distinct roots in $\mathbb K$ is diagonalisable: the polynomial has the shape the theorem asks for, and $\chi_A(A)=0$ by Theorem~\ref{thm:cayleyhamilton} below. For this special case one can avoid Cayley--Hamilton. Eigenvectors for distinct eigenvalues are independent, by the same $L_j(A)$ argument, and $n$ of them form a basis. A diagonalisable matrix also stays diagonalisable on every stable subspace, because the same polynomial kills the restriction. Exercise~\ref{ex:codiag} uses that fact.

\begin{remark}
Cayley--Hamilton says that $\chi_A(A)=0$. The proof below, Theorem~\ref{thm:cayleyhamilton}, uses the adjugate. Volume~3 proves the same identity again, as a fact about modules over $\mathbb K[X]$. The minimal polynomial is the sharper tool there. This chapter does not need it.
\end{remark}

\begin{example}[linear recurrences]
If $u_{n+2}=au_{n+1}+bu_n$, then $\binom{u_{n+1}}{u_n}=M^n\binom{u_1}{u_0}$ with $M=\begin{pmatrix}a&b\\1&0\end{pmatrix}$. The characteristic polynomial is $X^2-aX-b$. When $b\ne0$ and the polynomial has two distinct roots $r_1,r_2$, Theorem~\ref{thm:diag} says that $M$ is diagonalisable, and $u_n=\alpha r_1^n+\beta r_2^n$ for constants $\alpha,\beta$ fixed by $u_0$ and $u_1$.
\end{example}

If $b=0$, the recurrence is direct: $u_1$ is arbitrary and $u_n=a^{n-1}u_1$ for $n\ge2$, while $u_0$ is arbitrary. When $a=0$, every $u_n$ with $n\ge2$ is zero. This avoids treating $0^0$ as an initial condition. The double-root case is the other case the example leaves out. For $u_{n+2}=2u_{n+1}-u_n$, the matrix $\begin{pmatrix}2&-1\\1&0\end{pmatrix}$ has characteristic polynomial $(X-1)^2$ and is not diagonalisable: its only eigenspace is the line spanned by $(1,1)$. The solutions are $u_n=\alpha+\beta n$. The factor $n$ is what the missing eigenvector leaves behind. Volume~3 explains that factor by the Jordan form.

\section{Inner products and orthogonal projection}

An inner product measures lengths and right angles. On a real vector space it is a symmetric bilinear function $\langle v,w\rangle$ with $\langle v,v\rangle>0$ for $v\ne0$. The usual example on $\R^n$ is $\sum_iv_iw_i$. Put $\|v\|=\sqrt{\langle v,v\rangle}$. We say $v\perp w$ when their inner product is zero. A family is \emph{orthonormal} when $\langle e_i,e_j\rangle$ is $1$ for $i=j$ and $0$ otherwise. Expand the square: if $v\perp w$, then
\[
\|v+w\|^2=\|v\|^2+2\langle v,w\rangle+\|w\|^2=\|v\|^2+\|w\|^2.
\]
That is Pythagoras. Chapter~\ref{ch:inequalities} returns to the triangle inequality and to Cauchy--Schwarz as facts about averages. Here the expansion above is enough, and the case of two vectors in Theorem~\ref{thm:gram} will be Cauchy--Schwarz again, read as the area of a parallelogram.

For $z\in\C^n$ we also use the Euclidean length $\|z\|^2=\sum_i|z_i|^2$, a positive real number when $z\ne0$. A real square matrix is \emph{orthogonal} when $P^{\mathsf T}P=I$. Its columns are an orthonormal basis, and it preserves inner products because $\langle Px,Py\rangle=x^{\mathsf T}P^{\mathsf T}Py$.

The same geometry works on other spaces. On real polynomials of degree at most $d$, write $f=\sum a_kX^k$ and $g=\sum b_kX^k$, and set $\langle f,g\rangle=\sum a_kb_k$. Another inner product, given by an integral, appears in Exercise~\ref{ex:hilbert}.

\begin{theorem}[Gram--Schmidt and orthogonal projection]\label{thm:gramschmidt}
Let $F$ be a subspace of a finite-dimensional real inner product space $E$, with basis $v_1,\dots,v_p$.
\begin{enumerate}[label=(\alph*)]
\item (\textbf{Gram--Schmidt}) There is an orthonormal family $e_1,\dots,e_p$ with $\operatorname{Span}(e_1,\dots,e_k)=\operatorname{Span}(v_1,\dots,v_k)$ for each $k$, given by $w_k=v_k-\sum_{i<k}\langle v_k,e_i\rangle e_i$ and $e_k=w_k/\|w_k\|$.
\item (\textbf{projection}) Every $x\in E$ is uniquely $x=p(x)+r$ with $p(x)\in F$ and $r\perp F$; explicitly $p(x)=\sum_{i=1}^p\langle x,e_i\rangle e_i$. The map $p$ is linear with $p^2=p$, and $p(x)$ is the point of $F$ closest to $x$:
\[
\|x-p(x)\|<\|x-y\|\quad\text{for every }y\in F,\ y\ne p(x).
\]
We call $d(x,F)=\|x-p(x)\|$ the \emph{distance} from $x$ to $F$.
\end{enumerate}
\end{theorem}
\begin{proof}
(a) Proceed by induction on $k$. Suppose $e_1,\dots,e_{k-1}$ are orthonormal and span the same subspace as $v_1,\dots,v_{k-1}$. Then $w_k$ is $v_k$ minus a vector of that span, so $w_k\ne0$: the vector $v_k$ is not in the span of the earlier $v_i$, by independence. For $j<k$,
\[
\langle w_k,e_j\rangle=\langle v_k,e_j\rangle-\sum_{i<k}\langle v_k,e_i\rangle\langle e_i,e_j\rangle=\langle v_k,e_j\rangle-\langle v_k,e_j\rangle=0.
\]
Thus $e_k=w_k/\|w_k\|$ extends the orthonormal family. Also $\operatorname{Span}(e_1,\dots,e_k)\subseteq\operatorname{Span}(v_1,\dots,v_k)$. Both spaces have dimension $k$, so they are equal.

(b) Put $p(x)=\sum_{i\le p}\langle x,e_i\rangle e_i\in F$ and $r=x-p(x)$. For each $j$, one has $\langle r,e_j\rangle=\langle x,e_j\rangle-\langle x,e_j\rangle=0$, so $r\perp F$. For uniqueness, if $x=f+r=f'+r'$ with $f,f'\in F$ and $r,r'\perp F$, then $f-f'=r'-r$ lies in $F$ and is orthogonal to $F$, hence orthogonal to itself, so $f=f'$. Linearity and $p^2=p$ are read off the formula, since $p(e_i)=e_i$. For $y\in F$,
\[
x-y=(x-p(x))+(p(x)-y),
\]
and the two terms are orthogonal. Pythagoras gives
\[
\|x-y\|^2=\|x-p(x)\|^2+\|p(x)-y\|^2\ge\|x-p(x)\|^2,
\]
with equality only for $y=p(x)$.
\end{proof}

Gram--Schmidt shows that every subspace has an orthonormal basis, including the zero subspace, whose basis is empty. The projection formula is then $p(x)=0$. The proof uses only the inner product rules, so it applies to every finite-dimensional real inner product space. The right panel of Figure~\ref{fig:parallelogram} was this construction for two vectors in the plane. Figure~\ref{fig:gramschmidt} writes the arithmetic for $v_1=(3,1)$ and $v_2=(1,2)$. The closest-point property is the reason least squares works. The workshop at the end of the chapter is that property, with the data left in.

\begin{figure}[ht]
\centering
\fig{ch07-gramschmidt}
\caption{Left: the nearest point of a line $F$ to a point $x$ is the foot of the perpendicular, and any other point $y$ of the line is farther. Right: Gram--Schmidt for $v_1=(3,1)$ and $v_2=(1,2)$. The projection of $v_2$ on the line $\R v_1$ is $\langle v_2,e_1\rangle e_1=\frac12v_1$; the remainder $w_2=v_2-\frac12v_1=\bigl(-\frac12,\frac32\bigr)$ is perpendicular to $v_1$, and $\|w_2\|=\sqrt{10}/2$ is the distance from $v_2$ to the line. The area of the parallelogram is $\|v_1\|\,\|w_2\|=\sqrt{10}\cdot\sqrt{10}/2=5=\det\begin{pmatrix}3&1\\1&2\end{pmatrix}$, and the Gram determinant is $\det\begin{pmatrix}10&5\\5&5\end{pmatrix}=25=5^2$.}
\label{fig:gramschmidt}
\end{figure}

\section{The spectral theorem and quadratic forms}
\begin{think}
An invertible $2\times2$ symmetric matrix maps the unit circle to an ellipse (Figure~\ref{fig:ellipse}). An ellipse has two perpendicular axes. What should they be for the matrix, and what are their lengths?
\end{think}

\begin{theorem}[spectral theorem]\label{thm:spectral}
A real symmetric matrix $S$ has real eigenvalues and an orthonormal basis of eigenvectors: $S=PDP^{\mathsf T}$ with $P$ orthogonal and $D$ diagonal.
\end{theorem}
\begin{proof}
Let $\lambda\in\C$ be a root of $\chi_S$, and let $z\in\C^n\setminus\{0\}$ satisfy $Sz=\lambda z$. Then $\bar z^{\mathsf T}Sz=\lambda\|z\|^2$. This number is real: its conjugate is $z^{\mathsf T}S\bar z$, which is the same scalar transposed, because $S^{\mathsf T}=S$. So $\lambda\in\R$. The real and imaginary parts of $z$ both satisfy the real eigenvalue equation. At least one of them is nonzero. Normalising it gives a real unit eigenvector $u$.

Proceed by induction on $n$. The hyperplane $H=u^\perp$ is stable under $S$: if $x\perp u$, then $\langle Sx,u\rangle=\langle x,Su\rangle=\lambda\langle x,u\rangle=0$. The restriction of $S$ to $H$ is symmetric for the induced inner product, so $H$ has an orthonormal basis of eigenvectors. Add $u$.
\end{proof}

\begin{figure}[ht]
\centering
\fig{ch07-ellipse}
\caption{Left: $S=\begin{pmatrix}2&1\\1&2\end{pmatrix}$ maps the unit circle (grey) to an ellipse (charcoal). Each rust arrow joins a point $x$ of the circle to $Sx$. The axes of the ellipse are the orthonormal eigenvectors $u_1=\frac1{\sqrt2}(-1,1)$ and $u_2=\frac1{\sqrt2}(1,1)$, stretched by the eigenvalues $1$ and $3$, the extreme values of $x^{\mathsf T}Sx$ on the circle (Exercise~\ref{ex:rayleigh}). Right: $\det H_n$ for the Hilbert matrix of Exercise~\ref{ex:hilbert}, on a logarithmic scale. It is positive but falls from $\frac1{12}$ for $n=2$ to about $2.6\times10^{-78}$ for $n=12$.}
\label{fig:ellipse}
\end{figure}

The ellipse in Figure~\ref{fig:ellipse} has perpendicular axes because of this theorem. Cauchy's 1829 argument was about the same picture: a quadratic form in three variables, and the axes of the ellipsoid it defines. The probability chapter of this volume meets the same matrices again. A covariance matrix is positive semidefinite, because a variance is an average of squares, and the spectral theorem supplies an orthonormal basis of axes.

\begin{definitionx}[quadratic forms and positivity]
For a real symmetric matrix $S$, the function $q_S(x)=x^{\mathsf T}Sx$ is its \emph{quadratic form}. The matrix is \emph{positive semidefinite} if $q_S(x)\ge0$ for all $x$, and \emph{positive definite} if $q_S(x)>0$ for every $x\ne0$.
\end{definitionx}

The spectral theorem gives $q_S(Py)=\sum_i\lambda_i y_i^2$. Consequently $S$ is positive semidefinite exactly when every eigenvalue is nonnegative, and positive definite exactly when every eigenvalue is positive. One direction is the sum. The other is evaluation at a unit eigenvector. For example $\begin{pmatrix}1&1\\1&1\end{pmatrix}$ has quadratic form $(x_1+x_2)^2$, so it is semidefinite and not definite: the nonzero vector $(1,-1)$ gives zero.

\section{Determinants as volumes}

Back to the chessboard. The gap is a parallelogram, and its area was a $2\times2$ determinant. In higher dimension the volume is defined by repeating base times height. The projection theorem supplies the heights.

\begin{definitionx}[volume]
Let $v_1$, $\dots$, $v_p$ belong to $\R^n$. Define the successive heights by
\[
|\begin{gathered}
d_1=\|v_1\|,\\
d_k=d\left(v_k,\operatorname{Span}(v_1,\dots,v_{k-1})\right)\quad(k\ge2).
\end{gathered}
\]
The \emph{$p$-dimensional volume} of the parallelepiped $\{\sum t_iv_i:0\le t_i\le1\}$ is $\operatorname{vol}(v_1,\dots,v_p)=d_1d_2\cdots d_p$: base times height, repeated.
\end{definitionx}

\begin{theorem}[Gram determinant]\label{thm:gram}
Let $v_1,\dots,v_p\in\R^n$ and $G=\left(\langle v_i,v_j\rangle\right)_{1\le i,j\le p}$, the \emph{Gram matrix}. Then
\[
\det G=\operatorname{vol}(v_1,\dots,v_p)^2 ,
\]
so $\det G\ge0$, with equality if and only if $v_1,\dots,v_p$ are linearly dependent. If $p=n$ and $A$ is the matrix with columns $v_1,\dots,v_n$, then $G=A^{\mathsf T}A$ and
\[
|\det A|=\operatorname{vol}(v_1,\dots,v_n).
\]
\end{theorem}
\begin{proof}
If the $v_i$ are dependent, say $\sum c_iv_i=0$ with $c\ne0$, then $Gc=0$: the $i$th entry of $Gc$ is $\langle v_i,\sum c_jv_j\rangle$. So $\det G=0$. Some height $d_k$ is also zero, because the first $v_k$ that lies in the span of the earlier vectors has distance zero to that span. The volume is zero too.

Now suppose the $v_i$ are independent, and run Gram--Schmidt: $w_k=v_k-\sum_{i<k}\langle v_k,e_i\rangle e_i$. The subtracted vector is the projection of $v_k$ on $\operatorname{Span}(v_1,\dots,v_{k-1})$, so $\|w_k\|=d_k$. Replacing $v_k$ by $w_k$ replaces, in $G$, row $k$ by row $k$ minus a combination of the earlier rows, and column $k$ likewise. By Theorem~\ref{thm:det}(d), applied to rows and to columns, $\det G$ does not change. Do this for $k=2,\dots,p$ in order. Each step uses only the $w_i$ with $i<k$, and their span is the span of the $v_i$ with $i<k$. At the end, $G$ has become the Gram matrix of $w_1,\dots,w_p$, which are pairwise orthogonal. It is $\operatorname{diag}(\|w_1\|^2,\dots,\|w_p\|^2)$, with determinant $\prod d_k^2=\operatorname{vol}^2$.

For $p=n$, the $(i,j)$ entry of $A^{\mathsf T}A$ is $\langle v_i,v_j\rangle$, so $\det G=\det(A^{\mathsf T})\det A=(\det A)^2$ by Theorem~\ref{thm:det}(b).
\end{proof}

For $p=2$ the identity says $\|u\|^2\|v\|^2-\langle u,v\rangle^2\ge0$. That is Cauchy--Schwarz, with the equality case: a parallelogram has area zero exactly when its edges are parallel. For $p=n$ it says that the volume of the image of the unit cube under $A$ is $|\det A|$. Since $\det(AB)=\det A\det B$, the volume of the image of any parallelepiped is multiplied by $|\det A|$. Volume~7 proves the change-of-variables formula in which this factor, computed for the derivative, is called the Jacobian.

The same identity explains the Hilbert matrix of the epigraph. The matrix $H_n$ is the Gram matrix of $1,x,\dots,x^{n-1}$ for an inner product on polynomials (Exercise~\ref{ex:hilbert}), so $\det H_n$ is a squared volume and is positive. It is tiny because those monomials are nearly dependent for this inner product: the parallelepiped is almost flat. The right panel of Figure~\ref{fig:ellipse} plots the fall. Hadamard's inequality (Exercise~\ref{ex:hadamard}) is the same picture. A parallelepiped with given edge lengths has the largest volume when it is a box.

\section{The adjugate and Cayley--Hamilton}

Laplace expansion computes $\det A$ from the cofactors $c_{ij}$. If you expand along the wrong row, using the cofactors of one row and the entries of another, you get $0$. That remark is the inverse of a matrix, written in closed form.

\begin{proposition}[adjugate, Cramer's rule]\label{prop:adjugate}
Let $\operatorname{adj}A=(c_{ji})_{i,j}$ be the transpose of the matrix of cofactors. Then
\[
A\operatorname{adj}A=\operatorname{adj}A\;A=(\det A)\,I_n ,
\]
so if $\det A\ne0$, then $A^{-1}=\frac1{\det A}\operatorname{adj}A$. In that case the solution of $Ax=b$ is given by \textbf{Cramer's rule}: $x_j=\det A_j/\det A$, where $A_j$ is $A$ with column $j$ replaced by $b$.
\end{proposition}
\begin{proof}
The $(i,k)$ entry of $A\operatorname{adj}A$ is $\sum_ja_{ij}c_{kj}$. For $i=k$ this is the Laplace expansion of $\det A$ along row $i$. For $i\ne k$, it is the expansion along row $k$ of the matrix $A'$ obtained from $A$ by replacing row $k$ with a copy of row $i$. The cofactors of row $k$ do not use row $k$, so they are the same for $A$ and for $A'$, and the entries of row $k$ of $A'$ are the $a_{ij}$. But $A'$ has two equal rows, so $\det A'=0$. The same argument with columns gives $\operatorname{adj}A\;A=(\det A)I_n$.

For Cramer, $x=A^{-1}b=\frac1{\det A}(\operatorname{adj}A)\,b$ has $j$th entry $\frac1{\det A}\sum_kc_{kj}b_k$, and $\sum_kc_{kj}b_k$ is the expansion of $\det A_j$ along its column $j$. The entries of that column are the $b_k$, and its cofactors are those of $A$.
\end{proof}

For $n=2$ this is the formula in the worked example: $\operatorname{adj}\begin{pmatrix}a&b\\c&d\end{pmatrix}=\begin{pmatrix}d&-b\\-c&a\end{pmatrix}$. For large $n$, Cramer's rule is a theorem, not an algorithm. It costs $n+1$ determinants, while Gaussian elimination solves the system in about $n^3/3$ multiplications. What the rule is good for is structure. The entries of $A^{-1}$ are polynomials in the entries of $A$, divided by $\det A$. Exercise~\ref{ex:intinverse} uses that shape to decide which integer matrices have integer inverses. The same shape is the key to the next theorem. Cayley stated it in 1858, for $2\times2$ matrices with a remark on the general case. Frobenius proved it for every $n$ in 1878.

\portraits{cayley}{Arthur Cayley (1821--1895). Photograph by Herbert Rose Barraud, before 1883. Public domain, via Wikimedia Commons.}{hamilton}{William Rowan Hamilton (1805--1865). Portrait, artist unknown. Public domain, via Wikimedia Commons.}

\begin{theorem}[Cayley--Hamilton]\label{thm:cayleyhamilton}
Every square matrix satisfies its own characteristic equation: $\chi_A(A)=0$.
\end{theorem}

There is a one-line proof that is too good to be true. ``Since $\chi_A(X)=\det(XI-A)$, substitute $X=A$: then $\chi_A(A)=\det(AI-A)=\det0=0$.'' The symbol $X$ is being asked to be two things at once. Inside $\det(XI-A)$, it is a scalar on the diagonal. The expression $AI-A$ would have to mean ``the matrix $A$ in each diagonal slot'', which is a block matrix, not the zero matrix. The proof that works uses the adjugate of $XI-A$, a matrix whose entries are polynomials in $X$, and compares coefficients.

\begin{proof}
Write $\chi_A=X^n+c_{n-1}X^{n-1}+\dots+c_1X+c_0$. The cofactors of $XI_n-A$ are determinants of $(n-1)\times(n-1)$ matrices with entries in $\mathbb K[X]$ of degree at most $1$, so they are polynomials of degree at most $n-1$. Therefore
\[
\operatorname{adj}(XI_n-A)=B_0+XB_1+\dots+X^{n-1}B_{n-1}
\]
for some matrices $B_k\in M_n(\mathbb K)$, obtained by collecting the coefficient of $X^k$ in each entry. Proposition~\ref{prop:adjugate} is a polynomial identity, valid over any commutative ring, and here the ring is $\mathbb K[X]$. It gives
\[
(XI_n-A)\bigl(B_0+XB_1+\dots+X^{n-1}B_{n-1}\bigr)=\chi_A(X)\,I_n.
\]
Two matrices of polynomials are equal when their coefficients match. Compare powers of $X$:
\[
\begin{aligned}
X^n:&\quad B_{n-1}=I_n,\\
X^k\ (1\le k\le n-1):&\quad B_{k-1}-AB_k=c_kI_n,\\
X^0:&\quad -AB_0=c_0I_n.
\end{aligned}
\]
Multiply the equation of $X^k$ on the left by $A^k$, and add. The left side is
\[
A^nB_{n-1}+\sum_{k=1}^{n-1}\bigl(A^kB_{k-1}-A^{k+1}B_k\bigr)-AB_0.
\]
This is a telescoping sum. The term $A^kB_{k-1}$ cancels $-A^kB_{k-1}$ from the next equation down, and nothing remains. The right side is $A^n+\sum_{k<n}c_kA^k=\chi_A(A)$.
\end{proof}

The telescoping used only left multiplication by powers of $A$. It did not need the matrices $B_k$ to commute with $A$. They do commute with $A$, but that is a separate fact, and the proof does not need it.

\begin{examplex}[Cayley--Hamilton at work]
For $n=2$, the identity is $A^2-(\operatorname{tr}A)A+(\det A)I_2=0$. It is worth multiplying out once, by hand. For the Fibonacci matrix $M=\begin{pmatrix}1&1\\1&0\end{pmatrix}$, one has $M^2=M+I$, so every power of $M$ is $aM+bI$. Exercise~\ref{ex:fibch} turns that into identities for Fibonacci numbers. In general, $\chi_A(A)=0$ writes $A^n$ as a combination of $I,A,\dots,A^{n-1}$. Every higher power $A^m$, and every polynomial in $A$, therefore lies in the span of $I,A,\dots,A^{n-1}$. The algebra generated by one matrix has dimension at most $n$.

When $c_0=(-1)^n\det A\ne0$, divide $\chi_A(A)=0$ by $c_0$ and factor $A$ out of the rest:
\[
A^{-1}=-\frac1{c_0}\bigl(A^{n-1}+c_{n-1}A^{n-2}+\dots+c_1I_n\bigr).
\]
The inverse is a polynomial in the matrix. Exercise~\ref{ex:annul} is a special case.
\end{examplex}

\begin{beauty}
Cassini's identity is $\det(M^n)=(\det M)^n$: an arithmetic fact read off an area. A real symmetric matrix, and in particular a covariance matrix in the probability chapter, is diagonal in an orthonormal basis. The Google matrix of Exercise~\ref{ex:stochastic} need not be symmetric, and its interesting eigenvector is still a probability. Simultaneous diagonalisation (Exercise~\ref{ex:codiag}) is the same idea for two maps: if they commute and each is diagonalisable, they share one basis. The Hilbert matrix is positive definite, and still badly behaved in the monomial basis, so ordinary floating-point inversion can lose many digits. A small determinant alone does not diagnose that. Exercise~\ref{ex:cauchydet} gives the exact determinant, and the solution compares the greatest stretch with the least.
\end{beauty}

\section{Workshop: least squares, with the geometry left in}

Three measurements rarely lie on a perfect straight line. The question is which line is closest, and how large the miss is.

Legendre printed the method in 1805, in an appendix on the orbits of comets. Gauss gave his own account in 1809, in the \emph{Theoria motus}, while treating the orbit of Ceres. The priority is disputed. The calculation below is the projection theorem, applied to a column space.

\begin{examplex}[three points and an honest line]
Fit $y=a+bx$ to the points $(0,1)$, $(1,2)$ and $(2,2)$ by minimising the sum of squared vertical errors. Write
\[
A=\begin{pmatrix}1&0\\1&1\\1&2\end{pmatrix},\qquad
Y=\begin{pmatrix}1\\2\\2\end{pmatrix}.
\]
The columns are independent. The closest vector $A\theta$ to $Y$ in their span satisfies $A^{\mathsf T}(Y-A\theta)=0$, because the residual of an orthogonal projection is orthogonal to the subspace. So
\[
\begin{pmatrix}3&3\\3&5\end{pmatrix}\binom ab=\binom56.
\]
The right side is $A^{\mathsf T}Y$. Its first entry is $1+2+2=5$, and its second is $0\cdot1+1\cdot2+2\cdot2=6$. Subtract the first equation from the second: $2b=1$. Hence $b=\frac12$ and $a=\frac76$. The residual vector is $\bigl(-\frac16,\frac13,-\frac16\bigr)^{\mathsf T}$, and it is orthogonal to both columns of $A$. Its squared length is $\frac16$.
\end{examplex}

\begin{figure}[ht]\centering\fig{review07-fit}
\caption{The least-squares line $y=\frac76+\frac x2$ and the three observations. Dashed vertical segments show the residuals. The quantity minimised is the sum of the squares of those vertical lengths, not the perpendicular distances drawn in the plane of the page.}\end{figure}

\subsection*{Why is this actually a minimum?}
\textbf{Mission.} Prove uniqueness without differentiating a quadratic in two variables.

\textbf{Hint.} Orthogonality turns the squared norm of a sum into a sum of squared norms.

\textbf{Solution.} Let $\widehat\theta=\bigl(\frac76,\frac12\bigr)^{\mathsf T}$ and $r=Y-A\widehat\theta$. For every $h\in\R^2$, the residual $r$ is orthogonal to $Ah$, so
\[
\|Y-A(\widehat\theta+h)\|^2=\|r\|^2+\|Ah\|^2\ge\|r\|^2.
\]
Equality forces $Ah=0$, and independence of the columns gives $h=0$. This is both the minimum and the uniqueness. If the columns are dependent, the closest fitted vector is still unique, but more than one coefficient vector can produce it.

\begin{reminder}[small determinant, different question]
For $0<\varepsilon<1$, the matrix $\varepsilon I$ has tiny determinant $\varepsilon^n$ and condition number $1$ in the Euclidean norm: it scales every direction by the same factor. The matrix $\operatorname{diag}(M,1/M)$ has determinant $1$ and condition number $M^2$ for $M\ge1$. Conditioning compares the greatest stretch with the least. Their product is the determinant, up to sign. Invertibility and floating-point accuracy answer different questions.
\end{reminder}

\photograph{astronaut}{NASA astronaut Eileen Collins in an orange flight suit beside a helmet.}{Astronaut Eileen Collins. NASA, public domain; via scikit-image. The workshop turns three measured points into a matrix and a vector. The photograph is context for that habit of reading observations as coordinates. It is not a data set, and it is not an endorsement by the person shown.}

\section*{What this chapter establishes}

A matrix is a linear map written in a basis. A change of basis changes the entries and leaves the map alone. Rank--nullity says that the dimensions collapsed in the kernel are exactly the dimensions missing from the image. The determinant is the signed volume of the image of the unit cube, which is why Cassini's identity is an area of $1$. The nearest point of a subspace is the orthogonal projection. A real symmetric matrix is diagonal in an orthonormal basis.

The next chapter differentiates and integrates. Wallis's product for $\pi$, and the integral $\int_0^\infty e^{-x^2}\,\mathrm dx=\sqrt\pi/2$, come out of one family of integrals. The volume of a parallelepiped waits there only as a hint. When a change of variables multiplies volumes, the factor is a Jacobian, the absolute value of a determinant of derivatives. Volume~7 proves that formula.

Volume~2 measures a Gaussian integer $a+bi$ by the norm $a^2+b^2$, the determinant of the matrix in Exercise~\ref{ex:twobytwo}. Volume~3 proves Cayley--Hamilton again, from the structure of modules over $\mathbb K[X]$. Its appendices continue the count by symmetry: the Perron root of a positive matrix, the circulants a symmetry diagonalises, and the linear algebra method in combinatorics. Singular values and the Cauchy--Binet formula wait in Volume~8, beside the spectral theorem for operators. The Google matrix in the exercises is the first hint of that Perron root. An eigenvector can be a probability.

\section{Exercises}
\begin{exercise}[$\star$]\label{ex:fibmat}
Let $M=\begin{pmatrix}1&1\\1&0\end{pmatrix}$ and $(F_n)$ the Fibonacci numbers.
\begin{enumerate}[label=(\alph*)]
\item Show that $M^n=\begin{pmatrix}F_{n+1}&F_n\\F_n&F_{n-1}\end{pmatrix}$ for $n\ge1$, and deduce Cassini's identity $F_{n+1}F_{n-1}-F_n^2=(-1)^n$.
\item Diagonalise $M$ and deduce Binet's formula $F_n=\frac{\varphi^n-\psi^n}{\sqrt5}$, where $\varphi=\frac{1+\sqrt5}2$ and $\psi=\frac{1-\sqrt5}2$.
\end{enumerate}
\end{exercise}

\begin{exercise}[$\star$]\label{ex:nilp}
Let $A\in M_n(\C)$ be \emph{nilpotent}, meaning that $A^m=0$ for some integer $m\ge1$. Show that $\operatorname{tr}A=0$ and $\det(I_n+A)=1$.
\end{exercise}

\begin{exercise}[$\star$]\label{ex:mindet}
Compute the determinant of the $n\times n$ matrix $\left(\min(i,j)\right)_{1\le i,j\le n}$.
\end{exercise}

\begin{exercise}[$\star\star$]\label{ex:annul}
Let $A\in M_n(\R)$ with $A^2-3A+2I_n=0$. Show that $A$ is invertible, compute $A^{-1}$, and show that $A$ is diagonalisable. What are the possible eigenvalues?
\end{exercise}

\begin{exercise}[$\star\star$]\label{ex:abba}
Let $A,B\in M_n(\C)$. Show that $AB$ and $BA$ have the same characteristic polynomial.
\end{exercise}

\begin{exercise}[$\star\star$]\label{ex:rayleigh}
Let $S$ be real symmetric with eigenvalues $\lambda_1\le\dots\le\lambda_n$.
\begin{enumerate}[label=(\alph*)]
\item Show that $\max_{\|x\|=1}x^{\mathsf T}Sx=\lambda_n$ and $\min_{\|x\|=1}x^{\mathsf T}Sx=\lambda_1$.
\item Show that if $x^{\mathsf T}Sx\ge0$ for all $x$, there is a symmetric $R$ with $x^{\mathsf T}Rx\ge0$ for all $x$ and $R^2=S$.
\end{enumerate}
\end{exercise}

\begin{exercise}[$\star\star\star$]\label{ex:hilbert}
\emph{Bridge: Polynomials, integration and positive matrices.}
\emph{The Hilbert matrix (return after Chapter~\ref{ch:calculus}).} Let $H_n=\left(\frac1{i+j-1}\right)_{1\le i,j\le n}$. Show that $H_n$ is the Gram matrix of $1,x,\dots,x^{n-1}$ for the inner product
$\langle f,g\rangle=\int_0^1fg$, and deduce that $\det H_n>0$.
\end{exercise}

\begin{exercise}[$\star\star\star$]\label{ex:codiag}
Let $A,B\in M_n(\C)$ be diagonalisable with $AB=BA$. Show that there is a single invertible $P$ such that $P^{-1}AP$ and $P^{-1}BP$ are both diagonal.
\end{exercise}

\subsection*{Further exercises}

\begin{exercise}[$\star$]\label{ex:twobytwo}
\emph{Complex numbers as matrices.} To $z=a+ib$ associate $M(z)=\begin{pmatrix}a&-b\\b&a\end{pmatrix}$.
\begin{enumerate}[label=(\alph*)]
\item Show that $M(z+w)=M(z)+M(w)$, $M(zw)=M(z)M(w)$ and $\det M(z)=|z|^2$. Deduce $|zw|=|z||w|$ from Theorem~\ref{thm:det}(b).
\item Show that $M(e^{i\theta})$ is the rotation matrix $R_\theta$, and that the identity $R_\theta R_\varphi=R_{\theta+\varphi}$ is the identity $e^{i\theta}e^{i\varphi}=e^{i(\theta+\varphi)}$ of Chapter~\ref{ch:complex} in disguise: each is equivalent to the addition formulas for $\cos$ and $\sin$.
\item What are the eigenvalues of $M(z)$ over $\C$? Is $M(z)$ diagonalisable over $\R$?
\end{enumerate}
\end{exercise}

\begin{exercise}[$\star\star$]\label{ex:tracerank}
\emph{Trace and rank.}
\begin{enumerate}[label=(\alph*)]
\item Show that $\operatorname{tr}(AB)=\operatorname{tr}(BA)$ for $A\in M_{n,m}$, $B\in M_{m,n}$, and deduce that similar matrices have the same trace. Show that $AB-BA=I_n$ is impossible for $A,B\in M_n(\mathbb K)$ when $\mathbb K=\R$ or $\C$.
\item Show that the rank of the columns of a matrix equals the rank of its rows: $\operatorname{rk}A=\operatorname{rk}A^{\mathsf T}$. (Write $A=BC$ with $B$ having $r=\operatorname{rk}A$ columns.)
\end{enumerate}
\end{exercise}

\begin{exercise}[$\star\star$]\label{ex:hadamard}
\emph{Hadamard's inequality.} Let $A\in M_n(\R)$ have columns $v_1,\dots,v_n$.
\begin{enumerate}[label=(\alph*)]
\item Show that $|\det A|\le\|v_1\|\cdots\|v_n\|$, with equality iff the columns are pairwise orthogonal or one of them is $0$.
\item Deduce that if all entries of $A$ satisfy $|a_{ij}|\le1$, then $|\det A|\le n^{n/2}$. Find a $2\times2$ and a $4\times4$ matrix with entries $\pm1$ attaining the bound.
\end{enumerate}
\end{exercise}

\begin{exercise}[$\star\star$]\label{ex:sylvester}
\emph{Sylvester's determinant identity.} For $A\in M_{m,n}(\mathbb K)$ and $B\in M_{n,m}(\mathbb K)$, show that $\det(I_m+AB)=\det(I_n+BA)$. (Multiply the block matrix $\begin{pmatrix}I_m&-A\\B&I_n\end{pmatrix}$ on the left, then on the right, by $\begin{pmatrix}I_m&A\\0&I_n\end{pmatrix}$.) Deduce that the non-zero eigenvalues of $AB$ and $BA$ coincide, and compute $\det(I_n+uv^{\mathsf T})$ for column vectors $u,v$.
\end{exercise}

\begin{exercise}[$\star$]\label{ex:fibch}
\emph{Fibonacci by Cayley--Hamilton.} With $M$ as in Exercise~\ref{ex:fibmat}, use $M^2=M+I_2$ to show that $M^n=F_nM+F_{n-1}I_2$ for $n\ge1$, and deduce $F_{m+n}=F_mF_{n+1}+F_{m-1}F_n$ and $F_{2n}=F_n(F_{n+1}+F_{n-1})$. Which Fibonacci numbers divide $F_{2n}$? Which divide $F_{kn}$?
\end{exercise}

\begin{exercise}[$\star\star$]\label{ex:intinverse}
\emph{Integer matrices.} Let $A\in M_n(\Z)$.
\begin{enumerate}[label=(\alph*)]
\item Show that $A$ has an inverse with integer entries if and only if $\det A=\pm1$.
\item If $d=\det A\ne0$, show that $d\,\Z^n\subseteq A\Z^n$: every vector of $\Z^n$ divisible by $d$ is an integer combination of the columns of $A$.
\item Two bases of $\Z^2$ (pairs $u,v$ with $\Z u+\Z v=\Z^2$) span parallelograms of the same area. Which?
\end{enumerate}
\end{exercise}

\begin{exercise}[$\star\star\star$]\label{ex:stochastic}
\emph{A random surfer.} A matrix $G\in M_n(\R)$ is \emph{positive stochastic} if all its entries are $>0$ and each column sums to $1$. (The Google matrix of Figure~\ref{fig:pagerank} is one: $g_{ij}$ is the probability that a surfer on page $j$ moves to page $i$.)
\begin{enumerate}[label=(\alph*)]
\item Show that $1$ is an eigenvalue of $G$, and that every complex eigenvalue satisfies $|\lambda|\le1$.
\item Let $c=\min_{i,j}g_{ij}$ and let $\|x\|_1=\sum|x_i|$. Show that if $\sum_ix_i=0$ then $\|Gx\|_1\le(1-nc)\|x\|_1$.
\item Deduce that there is exactly one probability vector $\pi$ (non-negative entries summing to $1$) with $G\pi=\pi$, that all its entries are positive, and that for every probability vector $x_0$ the sequence $G^kx_0$ converges to $\pi$, with $\|G^kx_0-\pi\|_1\le2(1-nc)^k$.
\item Compute $\pi$ for the four-page web of Figure~\ref{fig:pagerank}.
\end{enumerate}
\end{exercise}

\begin{figure}[ht]
\centering
\fig{ch07-pagerank}
\caption{Left: a web of four pages and its links. Page $1$ links to $2$ and $3$; page $2$ to $3$; page $3$ to $1$; page $4$ to $1$ and $3$. Middle: the matrix $P$ of the surfer who clicks a random link ($p_{ij}=1/\text{out-degree}(j)$ if $j$ links to $i$), and the Google matrix $G=0.85P+0.15\cdot\frac14J$, where $J$ is the all-ones matrix: with probability $0.15$ the surfer jumps to a random page instead. Right: the entries of $G^kx_0$ for $x_0=(\frac14,\frac14,\frac14,\frac14)$; they converge geometrically to the PageRank vector $\pi$, with page $3$ ranked first and page $4$ last. Page~$4$ has no incoming link: its only visitors arrive by the random jump.}
\label{fig:pagerank}
\end{figure}

\begin{exercise}[$\star\star\star$]\label{ex:cauchydet}
\emph{Cauchy's determinant and the Hilbert matrix.} Let $x_1,\dots,x_n,y_1,\dots,y_n$ be scalars with all $x_i+y_j\ne0$.
\begin{enumerate}[label=(\alph*)]
\item Show that
\[
\det\left(\frac1{x_i+y_j}\right)_{1\le i,j\le n}=\frac{\prod_{1\le i<j\le n}(x_j-x_i)(y_j-y_i)}{\prod_{1\le i,j\le n}(x_i+y_j)}.
\]
\item Deduce $\det H_n=\dfrac{c_n^4}{c_{2n}}$ with $c_n=\prod_{k=1}^{n-1}k!$, and check $\det H_2=\frac1{12}$, $\det H_3=\frac1{2160}$.
\item Show that $\det H_n<4^{-(n-1)^2}$ for $n\ge2$: the determinant decays faster than any geometric sequence.
\end{enumerate}
\end{exercise}

\begin{exercise}[$\star\star$]\label{ex:bridge-translation}
\emph{Bridge: finite differences as a linear operator.} On $E=\R_d[X]$, $d\ge1$, define $S(P)(X)=P(X+1)$ and $\Delta=S-I$. Use the basis $b_k=\binom Xk$, $0\le k\le d$.
\begin{enumerate}[label=(\alph*)]
\item Find the matrices of $S$ and $\Delta$. Compute the characteristic polynomial and determinant of $S$.
\item Show that $\Delta^{d+1}=0$ but $\Delta^d\ne0$. Show that $S$ has only the eigenvalue $1$ and is not diagonalisable.
\item For every integer $m\ge0$, prove $S^m=\sum_{j=0}^d\binom mj\Delta^j$. Explain what this says about evaluating $P(X+m)$ from its finite differences.
\end{enumerate}
\end{exercise}

\begin{challenge}
\textbf{Three doors in the wall.} (1) \emph{Perron--Frobenius.} Exercise~\ref{ex:stochastic} assumed all entries positive. Show that for any matrix with positive entries (not necessarily stochastic) there is an eigenvalue $r>0$ which is a simple root of $\chi$, has an eigenvector with positive entries, and satisfies $|\lambda|<r$ for every other eigenvalue. Then ask what survives when some entries are $0$: the Google matrix adds $0.15J/n$ precisely to avoid this question.
(2) \emph{Hadamard's bound is sharp when it is.} A matrix with entries $\pm1$ and $|\det|=n^{n/2}$ is a \emph{Hadamard matrix}; show that it exists only for $n=1,2$ or $n\equiv0\pmod 4$. Whether it exists for \emph{every} multiple of $4$ is open; as of October 2026 the smallest undecided order is $2396$ (see the notes and sources).
(3) \emph{Every matrix is a product of two symmetric matrices} (Frobenius, 1910), over any field. Try $n=2$ by hand, then find why the general case is hard. Volume~3 gives the tools, through the rational canonical form.
\end{challenge}

\begin{exercise}[$\star$]\label{pb:tanroots}
\emph{The roots of $\tan x=x$.}
\begin{enumerate}[label=(\alph*)]
\item Show that for each $k\ge1$ the equation $\tan x=x$ has exactly one root $x_k$ in $\left]k\pi-\frac\pi2,k\pi+\frac\pi2\right[$, and that $k\pi<x_k<k\pi+\frac\pi2$.
\item With $a=(k+\frac12)\pi$, show that $x_k=a-\frac1a-\frac2{3a^3}+O(a^{-5})$.
\item Numerically, one finds $x_1=4.49341$ against $4.49381$ from (b). Explain why the approximation is already good for $k=1$.
\end{enumerate}
\end{exercise}

\section{Hints}

\begin{hint}{ex:fibmat}
(a) Induction, then take determinants. (b) The eigenvalues of $M$ are the roots of $X^2-X-1$.
\end{hint}
\begin{hint}{ex:nilp}
If $Av=\lambda v$ with $v\ne0$, then $A^kv=\lambda^kv$.
\end{hint}
\begin{hint}{ex:mindet}
Subtract each row from the next one, starting from the bottom.
\end{hint}
\begin{hint}{ex:annul}
$A(3I_n-A)=2I_n$, and $X^2-3X+2=(X-1)(X-2)$.
\end{hint}
\begin{hint}{ex:abba}
If $A$ is invertible, $BA=A^{-1}(AB)A$. Otherwise, replace $A$ by $A+tI_n$, and view both characteristic polynomials as polynomials in $t$.
\end{hint}
\begin{hint}{ex:rayleigh}
(a) Write $x=Py$ in an orthonormal eigenbasis: $x^{\mathsf T}Sx=\sum\lambda_iy_i^2$. (b) Take square roots of the eigenvalues.
\end{hint}
\begin{hint}{ex:hilbert}
$\frac1{i+j-1}=\int_0^1x^{i-1}x^{j-1}\dd x$. For $c\ne0$, what is $c^{\mathsf T}H_nc$?
\end{hint}
\begin{hint}{ex:codiag}
$B$ maps each eigenspace of $A$ into itself. Diagonalise $B$ inside each one.
\end{hint}
\begin{hint}{ex:twobytwo}
(a) Multiply out. (b) $e^{i\theta}=\cos\theta+i\sin\theta$. (c) $\chi=X^2-2aX+a^2+b^2$.
\end{hint}
\begin{hint}{ex:tracerank}
(a) Both traces are $\sum_{i,j}a_{ij}b_{ji}$. (b) If $b_1,\dots,b_r$ is a basis of the column space, each column of $A$ is a combination of them; write that as $A=BC$ and transpose.
\end{hint}
\begin{hint}{ex:hadamard}
(a) In Theorem~\ref{thm:gram}, $d_k\le\|v_k\|$. (b) Each column has norm $\le\sqrt n$. For $n=2$ take $H=\begin{pmatrix}1&1\\1&-1\end{pmatrix}$, for $n=4$ the block matrix $\begin{pmatrix}H&H\\H&-H\end{pmatrix}$.
\end{hint}
\begin{hint}{ex:sylvester}
The two products are block triangular, with blocks $I_m+AB$ and $I_n+BA$ on their diagonals. For the eigenvalues, replace $A$ by $-A/\lambda$.
\end{hint}
\begin{hint}{ex:fibch}
Induction for $M^n$; then compare the entries of $M^{m+n}=M^mM^n$ using Exercise~\ref{ex:fibmat}(a).
\end{hint}
\begin{hint}{ex:intinverse}
(a) $\det A\det A^{-1}=1$ with both factors integers; conversely use the adjugate. (b) $A\operatorname{adj}A=dI$. (c) The change-of-basis matrix is in $M_2(\Z)$ with integer inverse.
\end{hint}
\begin{hint}{ex:stochastic}
(a) The row vector $\mathbf 1=(1,\dots,1)$ satisfies $\mathbf 1G=\mathbf 1$; use $\det(G^{\mathsf T}-\lambda I)=\det(G-\lambda I)$. For $|\lambda|\le1$ consider a left eigenvector $y$ and the index where $|y_i|$ is largest. (b) Write each column as $c\mathbf 1+h_j$ with $h_j\ge0$. (c) The sequence $G^kx_0$ is Cauchy for $\|\cdot\|_1$; two fixed probability vectors have a difference of sum $0$.
\end{hint}
\begin{hint}{ex:cauchydet}
(a) Subtract the last column from each other column, factor $y_n-y_j$ out of column $j$ and $\frac1{x_i+y_n}$ out of row $i$; then do the same with rows, and expand along the last column, which has become $(0,\dots,0,1)^{\mathsf T}$. (b) $x_i=i$, $y_j=j-1$. (c) $c_{2n}/c_n^4=\prod_{k=1}^{n-1}\binom{2k}k^2(2k+1)$, and $\binom{2k}k\ge\frac{4^k}{2\sqrt k}$ (induction on $k$).
\end{hint}

\begin{hint}{ex:bridge-translation}
Chapter~\ref{ch:sums} gives $\Delta b_k=b_{k-1}$, with $\Delta b_0=0$. For diagonalisation, determine the dimension of the eigenspace, not just the roots of the characteristic polynomial.
\end{hint}

\begin{hint}{pb:tanroots}
(a) $(\tan x-x)'=\tan^2x$. (b) Put $x=a-\varepsilon$; then $\tan\varepsilon=\frac1{a-\varepsilon}$. Get $\varepsilon=O(\frac1a)$ first, then expand $\arctan$.
\end{hint}

\section{Solutions}

\begin{solution}{ex:fibmat}
\idea A two-term recurrence is a matrix power, so identities about $F_{n\pm1},F_n$ are identities about $M^n$; Cassini's left side looks like a $2\times2$ determinant, and $\det M^n=(\det M)^n$ is free. For Binet, diagonalising $M$ shows every entry of $M^n$ is a combination of $\varphi^n,\psi^n$; two initial values fix the coefficients.
\writeup (a) True for $n=1$ (with $F_0=0$), and $M^{n+1}=M^nM=\begin{pmatrix}F_{n+1}+F_n&F_{n+1}\\F_n+F_{n-1}&F_n\end{pmatrix}$ has the right form.
Taking determinants (Theorem~\ref{thm:det}(b)), $F_{n+1}F_{n-1}-F_n^2=(\det M)^n=(-1)^n$.

(b) $\chi_M=X^2-X-1$ has roots $\varphi\ne\psi$, so $M$ is diagonalisable (Theorem~\ref{thm:diag}): $M=PDP^{-1}$ and $M^n=PD^nP^{-1}$. Each entry of $M^n$ is therefore of the form
$\alpha\varphi^n+\beta\psi^n$; for $F_n$, the values $F_0=0$, $F_1=1$ give $\alpha+\beta=0$, $\alpha\varphi+\beta\psi=1$, so $\alpha=-\beta=\frac1{\varphi-\psi}=\frac1{\sqrt5}$.
\end{solution}

\begin{solution}{ex:nilp}
\idea Trace and determinant are read off the eigenvalues (Proposition~\ref{prop:charpoly}), so find the eigenvalues of a nilpotent matrix: apply $A^k=0$ to an eigenvector. Then $I_n+A$ has its eigenvalues shifted by $1$.
\writeup If $A^k=0$ and $Av=\lambda v$ with $v\ne0$, then $0=A^kv=\lambda^kv$, so $\lambda=0$. Over $\C$, $\chi_A$ splits, and all its roots are $0$: $\chi_A=X^n$. Hence
$\operatorname{tr}A=0$, and $\chi_{I+A}(X)=\det((X-1)I-A)=\chi_A(X-1)=(X-1)^n$, whose constant term $(-1)^n$ equals $(-1)^n\det(I_n+A)$: so $\det(I_n+A)=1$.
\end{solution}

\begin{solution}{ex:mindet}
\idea The entries $\min(i,j)$ grow by $1$ from one row to the next, but only in the columns $j\ge i$. So the difference of consecutive rows is a row of $0$s then $1$s, and row operations (Theorem~\ref{thm:det}(d)) should make the matrix triangular. Work from the bottom up so that each subtraction uses an original row.
\writeup For $i=n,n-1,\dots,2$ in that order, subtract row $i-1$ from row $i$. The new row $i$ has entries $\min(i,j)-\min(i-1,j)$, which is $1$ for $j\ge i$ and $0$ for $j<i$.
The matrix becomes upper triangular with ones on the diagonal: the determinant is $1$.
\end{solution}

\begin{solution}{ex:annul}
\idea An annihilating polynomial with non-zero constant term gives the inverse: isolate the constant and factor $A$ out of the rest. The polynomial factors as $(X-1)(X-2)$, split with simple roots, which is exactly the criterion of Theorem~\ref{thm:diag}.
\writeup $A(3I_n-A)=2I_n$, so $A$ is invertible with $A^{-1}=\frac12(3I_n-A)$. The polynomial $X^2-3X+2=(X-1)(X-2)$ is split with simple roots and cancels $A$,
so $A$ is diagonalisable, with eigenvalues in $\{1,2\}$ (if $Av=\lambda v$, $v\ne0$, then $0=(A^2-3A+2I)v=(\lambda^2-3\lambda+2)v$). Each value can occur: $I_n$ has only $1$, and $2I_n$ only $2$. For $n\ge2$, a diagonal matrix with both $1$ and $2$ has both eigenvalues.
\end{solution}

\begin{solution}{ex:abba}
\idea When $A$ is invertible, $BA$ is a conjugate of $AB$, and similar matrices have the same $\chi$. The general case is a \emph{density argument}: perturb $A$ to $A+tI$, invertible for all but finitely many $t$, and use that a polynomial identity in $t$ true at infinitely many points is true everywhere.
\writeup If $A$ is invertible, $BA=A^{-1}(AB)A$ is similar to $AB$, so they have the same characteristic polynomial (Proposition~\ref{prop:charpoly}). In general, $A+tI_n$ is invertible for all but finitely many $t\in\C$
(those $-t$ that are eigenvalues of $A$). For such $t$ and every $x$,
\[
\begin{aligned}&\det\left(xI_n-(A+tI_n)B\right)\eqbr=\det\left(xI_n-B(A+tI_n)\right).\end{aligned}
\]
For fixed $x$, both sides are polynomials in $t$ that agree at infinitely many points, so they are equal (Proposition~\ref{prop:rootsfactor}), in particular at $t=0$.
\end{solution}

\begin{solution}{ex:rayleigh}
\idea Questions about a symmetric matrix become questions about a diagonal one in an orthonormal eigenbasis (Theorem~\ref{thm:spectral}), and orthogonal changes of basis preserve unit vectors. Then $x^{\mathsf T}Sx$ is a weighted average $\sum\lambda_iy_i^2$ with weights $y_i^2$ summing to $1$. For (b), take square roots of the eigenvalues, which is possible since they are $\ge0$ by (a).
\writeup (a) Write $S=PDP^{\mathsf T}$ and $y=P^{\mathsf T}x$; $P$ is orthogonal so $\|y\|=\|x\|=1$, and $x^{\mathsf T}Sx=y^{\mathsf T}Dy=\sum\lambda_iy_i^2$. This lies between $\lambda_1\sum y_i^2=\lambda_1$
and $\lambda_n$, and the bounds are attained at unit eigenvectors.

(b) By (a) all $\lambda_i\ge0$. Let $R=P\sqrt DP^{\mathsf T}$ where $\sqrt D=\operatorname{diag}(\sqrt{\lambda_i})$. Then $R$ is symmetric, $R^2=PDP^{\mathsf T}=S$, and its eigenvalues
$\sqrt{\lambda_i}\ge0$ give $x^{\mathsf T}Rx\ge0$ by (a).
\end{solution}

\begin{solution}{ex:hilbert}
\idea The number $\frac1{i+j-1}$ is the integral $\int_0^1x^{i-1}x^{j-1}\dd x$, and seeing it that way turns the matrix into a Gram matrix. Gram matrices of independent vectors are positive definite, because $c^{\mathsf T}Gc$ is a squared norm; positive definite symmetric matrices have positive eigenvalues, hence positive determinant.
\writeup On the finite-dimensional space of real polynomials of degree at most $n-1$, put $\langle f,g\rangle=\int_0^1fg$. Bilinearity and symmetry follow from the integral rules. If $f\ne0$, some point of $[0,1]$ has $f(x)\ne0$ by the polynomial root bound; continuity gives a subinterval of positive length on which $f^2$ is bounded below by a positive number. Thus $\int_0^1f^2>0$, proving positive definiteness. For this inner product,
\[
\langle x^{i-1},x^{j-1}\rangle=\int_0^1x^{i+j-2}\dd x=\frac1{i+j-1}.
\]
For a Gram matrix $G$ with entries $\langle e_i,e_j\rangle$, of linearly independent vectors and $c\in\R^n\setminus\{0\}$,
$c^{\mathsf T}Gc=\left\|\sum c_ie_i\right\|^2>0$, because $\sum c_ie_i\ne0$. So $G$ is symmetric with all eigenvalues positive (by Exercise~\ref{ex:rayleigh}(a), the least one is $\min_{\|c\|=1}c^{\mathsf T}Gc>0$), and $\det G$, their product, is positive.
The monomials $1,x,\dots,x^{n-1}$ are linearly independent (a non-zero polynomial has finitely many roots), so $\det H_n>0$.
\end{solution}

\begin{solution}{ex:codiag}
\idea Commuting matrices preserve each other's eigenspaces. So decompose the space into eigenspaces of $A$ and diagonalise $B$ inside each one. That $B$ stays diagonalisable on a stable subspace follows from the polynomial criterion: a polynomial that kills $B$ kills its restriction.
\writeup Let $\lambda_1,\dots,\lambda_r$ be the distinct eigenvalues of $A$ and $E_k=\ker(A-\lambda_kI_n)$; since $A$ is diagonalisable, $\C^n=E_1\oplus\dots\oplus E_r$ (Theorem~\ref{thm:diag}).
Each $E_k$ is stable under $B$: if $Av=\lambda_kv$ then $A(Bv)=BAv=\lambda_kBv$. A polynomial with simple roots that cancels $B$ also cancels its restriction $B_k$ to $E_k$,
so $B_k$ is diagonalisable: $E_k$ has a basis of eigenvectors of $B$, which are automatically eigenvectors of $A$ (for $\lambda_k$). Gathering these bases gives a basis of $\C^n$ made of common eigenvectors; let $P$ have them as columns.
\end{solution}


\begin{solution}{ex:twobytwo}
\idea Everything is a direct computation; the point is to notice that the computation of $M(z)M(w)$ \emph{is} the computation of $(a+ib)(c+id)$, with the same four products in the same places. So the complex numbers are a copy of a family of $2\times2$ matrices, and matrix facts (determinants multiply) become facts about complex numbers (moduli multiply).
\writeup (a) Additivity is clear. For $z=a+ib$, $w=c+id$,
\[
M(z)M(w)=\begin{pmatrix}ac-bd&-(ad+bc)\\ad+bc&ac-bd\end{pmatrix}=M\big((ac-bd)+i(ad+bc)\big)=M(zw),
\]
and $\det M(z)=a^2+b^2=|z|^2$. Hence $|zw|^2=\det M(zw)=\det M(z)\det M(w)=|z|^2|w|^2$.

(b) Since $e^{i\theta}=\cos\theta+i\sin\theta$,
\[
M(e^{i\theta})=\begin{pmatrix}\cos\theta&-\sin\theta\\\sin\theta&\cos\theta\end{pmatrix}=R_\theta.
\] Since $M$ is injective and multiplicative, $R_\theta R_\varphi=R_{\theta+\varphi}$ holds iff $M(e^{i\theta}e^{i\varphi})=M(e^{i(\theta+\varphi)})$ iff $e^{i\theta}e^{i\varphi}=e^{i(\theta+\varphi)}$. Reading the first column of either identity gives $\cos(\theta+\varphi)=\cos\theta\cos\varphi-\sin\theta\sin\varphi$ and $\sin(\theta+\varphi)=\sin\theta\cos\varphi+\cos\theta\sin\varphi$, and conversely these two formulas give the identity.

(c) $\chi_{M(z)}=(X-a)^2+b^2$, with roots $a\pm ib=z,\bar z$; the eigenvectors are $(1,-i)$ for $z$ and $(1,i)$ for $\bar z$ (check: $M(z)\binom1{-i}=\binom{a+ib}{b-ia}=z\binom1{-i}$). Over $\R$, if $b\ne0$ there is no real eigenvalue at all, so $M(z)$ is not diagonalisable; if $b=0$ it is the scalar matrix $aI_2$. A rotation of the plane by an angle that is not a multiple of $\pi$ fixes no direction, which is the geometric content.
\end{solution}

\begin{solution}{ex:tracerank}
\idea (a) Write the trace as a double sum and swap the order. (b) The rank is the number of columns one needs to generate all of them; writing ``each column is a combination of $r$ chosen columns'' as a matrix product $A=BC$ and transposing turns a statement about columns into one about rows.
\writeup (a) $\operatorname{tr}(AB)=\sum_i\sum_ja_{ij}b_{ji}=\sum_j\sum_ib_{ji}a_{ij}=\operatorname{tr}(BA)$. Then $\operatorname{tr}(P^{-1}AP)=\operatorname{tr}(APP^{-1})=\operatorname{tr}A$, so the trace of a linear map is defined. If $AB-BA=I_n$, taking traces gives $0=n$, impossible in $\R$ or $\C$. (In infinite dimension the identity is possible: on polynomials, $(\frac{d}{dx}\circ x-x\circ\frac{d}{dx})P=P$. This is the canonical commutation relation: the trace argument proves that this exact relation has no finite-dimensional real or complex representation.)

(b) Let $r=\operatorname{rk}A$ be the dimension of the column space and $b_1,\dots,b_r$ a basis of it; let $B\in M_{n,r}$ have these as columns. Each column $a_j$ of $A$ is $\sum_kc_{kj}b_k$, which says $A=BC$ with $C=(c_{kj})\in M_{r,m}$. Transposing, $A^{\mathsf T}=C^{\mathsf T}B^{\mathsf T}$, whose columns are combinations of the $r$ columns of $C^{\mathsf T}$: so $\operatorname{rk}A^{\mathsf T}\le r=\operatorname{rk}A$. Applying this to $A^{\mathsf T}$ gives the reverse inequality.
\end{solution}

\begin{solution}{ex:hadamard}
\idea Theorem~\ref{thm:gram} writes $|\det A|$ as a product of heights $d_k$, and a height is at most the length of the edge: that is the whole proof, and equality means every edge is its own height, that is, perpendicular to the previous ones. For (b), a vector with $n$ entries of size $\le1$ has length $\le\sqrt n$.
\writeup (a) If the columns are dependent, $\det A=0\le\prod\|v_k\|$, with equality iff some $v_k=0$ (if all $v_k\ne0$ the product is positive). Otherwise, by Theorem~\ref{thm:gram}, $|\det A|=\prod_kd_k$ with $d_k=\|v_k-p_k(v_k)\|$, $p_k$ the projection on $F_k=\operatorname{Span}(v_1,\dots,v_{k-1})$. Pythagoras gives $\|v_k\|^2=\|p_k(v_k)\|^2+d_k^2$, so $d_k\le\|v_k\|$, and equality in the product (all factors positive) forces $p_k(v_k)=0$ for every $k$, that is $v_k\perp F_k$: the columns are pairwise orthogonal. Conversely, pairwise orthogonal non-zero columns have $p_k(v_k)=0$ and give equality.

(b) $\|v_k\|^2=\sum_ia_{ik}^2\le n$, so $|\det A|\le n^{n/2}$. Equality needs entries $\pm1$ and orthogonal columns. $H=\begin{pmatrix}1&1\\1&-1\end{pmatrix}$ has $|\det H|=2=2^{2/2}$, and $H_4=\begin{pmatrix}H&H\\H&-H\end{pmatrix}$ has orthogonal columns of norm $2$, so $|\det H_4|=2^4=16=4^{4/2}$. (The construction $H\mapsto\begin{pmatrix}H&H\\H&-H\end{pmatrix}$ doubles the order, so Hadamard matrices exist for every power of $2$; this is Sylvester's construction of 1867.)
\end{solution}

\begin{solution}{ex:sylvester}
\idea A block matrix with an identity block can be made block triangular by block row or column operations, and two different ways of doing it expose $I_m+AB$ and $I_n+BA$ as the only non-trivial diagonal blocks. The eigenvalue statement is the identity with $A$ scaled by $-1/\lambda$, and $\det(I+uv^{\mathsf T})$ is the case $m=n$, $n=1$ in reverse: a $1\times1$ determinant.
\writeup Let $M=\begin{pmatrix}I_m&-A\\B&I_n\end{pmatrix}$ and $T=\begin{pmatrix}I_m&A\\0&I_n\end{pmatrix}$, with $\det T=1$ (triangular). Block multiplication gives
\[
TM=\begin{pmatrix}I_m+AB&0\\B&I_n\end{pmatrix},\qquad MT=\begin{pmatrix}I_m&0\\B&I_n+BA\end{pmatrix}.
\]
For a block lower triangular matrix $\begin{pmatrix}X&0\\Y&Z\end{pmatrix}$ the determinant is $\det X\det Z$: indeed the matrix equals
\[
\begin{pmatrix}X&0\\0&I\end{pmatrix}\begin{pmatrix}I&0\\Y&I\end{pmatrix}\begin{pmatrix}I&0\\0&Z\end{pmatrix},
\]
the middle factor is triangular with determinant $1$, and $\det\begin{pmatrix}X&0\\0&I\end{pmatrix}=\det X$ because, as a function of the columns of $X$, it is multilinear, alternating and equal to $1$ at $X=I$ (Theorem~\ref{thm:det}(a)); likewise for $Z$. So $\det M=\det(I_m+AB)=\det(I_n+BA)$.

For $\lambda\ne0$, applying this to $-A/\lambda$ and $B$:
\[
\det(\lambda I_m-AB)=\lambda^m\det\left(I_m-\tfrac1\lambda AB\right)=\lambda^m\det\left(I_n-\tfrac1\lambda BA\right)=\lambda^{m-n}\det(\lambda I_n-BA),
\]
so $\lambda^n\chi_{AB}(\lambda)=\lambda^m\chi_{BA}(\lambda)$ for all $\lambda\ne0$, hence as polynomials: $\chi_{AB}$ and $\chi_{BA}$ differ by a power of $X$, and their non-zero roots coincide with multiplicities. (For $m=n$ this is Exercise~\ref{ex:abba} again.) Finally, with $A=u$ and $B=v^{\mathsf T}$ ($n=1$ on the small side), $\det(I_n+uv^{\mathsf T})=\det(1+v^{\mathsf T}u)=1+v^{\mathsf T}u=1+\sum_i u_iv_i$.
\end{solution}

\begin{solution}{ex:fibch}
\idea Cayley--Hamilton reduces every power of $M$ to the plane spanned by $I$ and $M$, and the coefficients are forced to be Fibonacci numbers by the recurrence. Then $M^{m+n}=M^mM^n$ is a product of two such expressions, and comparing coefficients gives the addition formula for free. Everything in the exercise is this one idea.
\writeup For $n=1$, $M=F_1M+F_0I$. If $M^n=F_nM+F_{n-1}I$, then $M^{n+1}=F_nM^2+F_{n-1}M=F_n(M+I)+F_{n-1}M=F_{n+1}M+F_nI$. Now
\[
\begin{aligned}
M^{m+n}&=(F_mM+F_{m-1}I)(F_nM+F_{n-1}I)\\
&=F_mF_n(M+I)+(F_mF_{n-1}+F_{m-1}F_n)M+F_{m-1}F_{n-1}I\\
&=(F_mF_{n+1}+F_{m-1}F_n)M+(F_mF_n+F_{m-1}F_{n-1})I ,
\end{aligned}
\]
and also $M^{m+n}=F_{m+n}M+F_{m+n-1}I$. Since $M$ is not a scalar matrix, $I$ and $M$ are linearly independent and the coefficients match: $F_{m+n}=F_mF_{n+1}+F_{m-1}F_n$. With $m=n$, $F_{2n}=F_n(F_{n+1}+F_{n-1})$; the second factor is the Lucas number $L_n$. So $F_n$ divides $F_{2n}$, and by induction on $k$, $F_n$ divides $F_{kn}$: $F_{(k+1)n}=F_{kn}F_{n+1}+F_{kn-1}F_n$ with both terms divisible by $F_n$. (Exercise~\ref{ex:fibgcd} in Chapter~\ref{ch:sums} proves the stronger identity $\gcd(F_m,F_n)=F_{\gcd(m,n)}$.)
\end{solution}

\begin{solution}{ex:intinverse}
\idea The adjugate formula shows that $A^{-1}$ has entries in $\frac1{\det A}\Z$, so the only obstacle to an integer inverse is the determinant; and $\det A\cdot\det A^{-1}=1$ with two integers leaves only $\pm1$. For (c), a change of basis of $\Z^2$ is an integer matrix with integer inverse, so its determinant is $\pm1$ and it preserves areas.
\writeup (a) If $A^{-1}\in M_n(\Z)$, then $\det A\cdot\det A^{-1}=1$ with both factors in $\Z$, so $\det A=\pm1$. Conversely, if $\det A=\pm1$, then $A^{-1}=\pm\operatorname{adj}A$ (Proposition~\ref{prop:adjugate}), and the cofactors are determinants of integer matrices, hence integers.

(b) For $w\in\Z^n$, $dw=A\operatorname{adj}(A)w$ and $\operatorname{adj}(A)w\in\Z^n$: so $dw$ is an integer combination of the columns of $A$.

(c) If $(u,v)$ and $(u',v')$ are bases of $\Z^2$, write $u'=au+cv$, $v'=bu+dv$ with $P=\begin{pmatrix}a&b\\c&d\end{pmatrix}\in M_2(\Z)$; expressing $u,v$ in terms of $u',v'$ gives an integer matrix which is $P^{-1}$, so $\det P=\pm1$ by (a), and $|\det(u',v')|=|\det P|\,|\det(u,v)|=|\det(u,v)|$. Comparing with the basis $(e_1,e_2)$: every basis of $\Z^2$ spans a parallelogram of area exactly $1$. Conversely, if $u,v\in\Z^2$ span area $1$, the matrix $(u\ v)$ has determinant $\pm1$, so by (a) its inverse is integral, and $e_1,e_2$ are integer combinations of $u,v$: it is a basis. The same inverse shows that the parallelogram contains no other lattice point. If $p=su+tv$ lies in $\Z^2$ with $0\le s,t\le 1$, then $\binom st=P^{-1}p$ has integer entries, so $s$ and $t$ are $0$ or $1$: the corners.
\end{solution}

\begin{solution}{ex:stochastic}
\idea Column sums equal to $1$ say that the row vector of ones is a left eigenvector for $1$. The contraction in (b) is the heart: subtracting from each column the constant vector $c\mathbf 1$ leaves non-negative columns of sum $1-nc$, and the constant part is invisible to vectors of sum $0$. After that, (c) is the standard route from ``contraction'' to ``unique fixed point, reached geometrically'': Cauchy sequences converge, and two fixed points would have to be at distance $0$.
\writeup (a) $\mathbf 1^{\mathsf T}G=\mathbf 1^{\mathsf T}$ because each column of $G$ sums to $1$; so $G^{\mathsf T}\mathbf 1=\mathbf 1$, $1$ is an eigenvalue of $G^{\mathsf T}$, and $\chi_{G^{\mathsf T}}=\chi_G$ (Theorem~\ref{thm:det}(b)), so of $G$. If $\lambda\in\C$ is an eigenvalue of $G$, it is one of $G^{\mathsf T}$: there is $y\ne0$ with $y^{\mathsf T}G=\lambda y^{\mathsf T}$, that is $\lambda y_j=\sum_iy_ig_{ij}$ for every $j$. Choose $j$ with $|y_j|$ maximal; then $|\lambda||y_j|\le\sum_i|y_i|g_{ij}\le|y_j|\sum_ig_{ij}=|y_j|$, so $|\lambda|\le1$.

(b) Write column $j$ of $G$ as $g_j=c\mathbf 1+h_j$ with $h_j\ge0$ entrywise and $\sum_ih_{ij}=1-nc$. For $x$ with $\sum_jx_j=0$,
\[
Gx=\sum_jx_jg_j=c\Big(\sum_jx_j\Big)\mathbf 1+\sum_jx_jh_j=\sum_jx_jh_j ,
\]
so $\|Gx\|_1\le\sum_j|x_j|\,\|h_j\|_1=(1-nc)\|x\|_1$. (Note $nc\le1$, since a column has $n$ entries $\ge c$ summing to $1$.) Put $q=1-nc<1$.

(c) $G$ maps probability vectors to probability vectors (non-negative entries, and $\mathbf 1^{\mathsf T}Gx=\mathbf 1^{\mathsf T}x=1$). Let $x_k=G^kx_0$. Then $x_{k+1}-x_k=G^k(x_1-x_0)$ with $x_1-x_0$ of sum $0$, so $\|x_{k+1}-x_k\|_1\le q^k\|x_1-x_0\|_1\le2q^k$ by (b) applied $k$ times (each $G^jx$ keeps sum $0$). Hence $\|x_l-x_k\|_1\le2\sum_{j\ge k}q^j=\frac{2q^k}{1-q}$ for $l>k$: the sequence is Cauchy in each coordinate, so it converges (Chapter~\ref{ch:sequences}) to some $\pi$, which is a probability vector (limits of non-negative numbers summing to $1$), and $G\pi=\lim Gx_k=\lim x_{k+1}=\pi$. If $\pi'$ is another fixed probability vector, $\pi-\pi'$ has sum $0$ and $\|\pi-\pi'\|_1=\|G(\pi-\pi')\|_1\le q\|\pi-\pi'\|_1$, forcing $\pi=\pi'$. Each entry is positive: $\pi_i=\sum_jg_{ij}\pi_j\ge c\sum_j\pi_j=c>0$. Finally $x_k-\pi=G^k(x_0-\pi)$ with $x_0-\pi$ of sum $0$ and $\|x_0-\pi\|_1\le2$, so $\|x_k-\pi\|_1\le2q^k$.

(d) Here $c=0.15/4=0.0375$ and $q=0.85$. The last row of $G$ is constant, equal to $0.0375$, so $\pi_4=0.0375\sum_j\pi_j=0.0375$ exactly: page~$4$ gets precisely the random-jump traffic. Solving $(G-I)\pi=0$ with $\sum\pi_j=1$ (three unknowns left) gives
\[
\pi=\left(\tfrac{2687}{7076},\ \tfrac{56293}{283040},\ \tfrac{108653}{283040},\ \tfrac{3}{80}\right)\approx(0.3797,\ 0.1989,\ 0.3839,\ 0.0375):
\]
page~$3$ ranks first, just ahead of page~$1$, which it links to. Figure~\ref{fig:pagerank} shows the power iteration; the actual convergence factor is $0.60$, the modulus of the other eigenvalues ($\chi_G=X(X-1)(X^2+0.85X+0.36125)$), better than the guaranteed $q=0.85$.
\end{solution}

\begin{solution}{ex:cauchydet}
\idea The entries $\frac1{x_i+y_j}$ have a product structure that row and column operations reveal: the difference of two such entries factors as $\frac{(y_n-y_j)}{(x_i+y_j)(x_i+y_n)}$. Clearing one column of $1$s and one row of $0$s costs a known product of factors and leaves the same determinant one size smaller: induction. The Hilbert matrix is the case $x_i=i$, $y_j=j-1$, and then every product is a product of factorials.
\writeup (a) Let $D_n$ be the determinant. Subtract column $n$ from each column $j<n$: entry $(i,j)$ becomes
\[
\frac1{x_i+y_j}-\frac1{x_i+y_n}=\frac{y_n-y_j}{(x_i+y_j)(x_i+y_n)}.
\]
Factor $y_n-y_j$ out of column $j$ ($j<n$) and $\frac1{x_i+y_n}$ out of row $i$ (all $i$): the matrix now has entries $\frac1{x_i+y_j}$ for $j<n$ and $1$ in column $n$. Subtract row $n$ from each row $i<n$: entry $(i,j)$, $j<n$, becomes $\frac{x_n-x_i}{(x_i+y_j)(x_n+y_j)}$, and column $n$ becomes $(0,\dots,0,1)^{\mathsf T}$. Factor $x_n-x_i$ out of row $i<n$ and $\frac1{x_n+y_j}$ out of column $j<n$. Expanding along the last column,
\[
D_n=\frac{\prod_{j<n}(y_n-y_j)\prod_{i<n}(x_n-x_i)}{\prod_{i\le n}(x_i+y_n)\prod_{j<n}(x_n+y_j)}\;D_{n-1},
\]
and $D_1=\frac1{x_1+y_1}$. Multiplying these relations for $n,n-1,\dots,2$ gives the formula: the numerators assemble $\prod_{i<j}(x_j-x_i)(y_j-y_i)$ and the denominators every $x_i+y_j$ exactly once.

(b) With $x_i=i$, $y_j=j-1$, $x_i+y_j=i+j-1$. For fixed $j$, $\prod_{i<j}(j-i)=(j-1)!$, so $\prod_{i<j}(x_j-x_i)=\prod_{j=1}^n(j-1)!=c_n$, and the same for the $y$'s: the numerator is $c_n^2$. The denominator is $\prod_{i=1}^n\prod_{j=1}^n(i+j-1)=\prod_{i=1}^n\frac{(i+n-1)!}{(i-1)!}=\frac{\prod_{m=n}^{2n-1}m!}{c_n}=\frac{c_{2n}}{c_n^2}$. So $\det H_n=c_n^4/c_{2n}$. Check: $c_2=1$, $c_4=1!\,2!\,3!=12$, $\det H_2=\frac1{12}$; $c_3=2$, $c_6=1!\,2!\,3!\,4!\,5!=34560$, $\det H_3=\frac{16}{34560}=\frac1{2160}$.

(c) Grouping $m!$ for $m=2k$ and $2k+1$, $c_{2n}=\prod_{k=1}^{n-1}(2k)!(2k+1)!$ (the factor $1!$ is $1$), so
\[
\frac{c_{2n}}{c_n^4}=\prod_{k=1}^{n-1}\frac{(2k)!\,(2k+1)!}{(k!)^4}=\prod_{k=1}^{n-1}\binom{2k}k^2(2k+1).
\]
Now $\binom{2k}k\ge\frac{4^k}{2\sqrt k}$: true for $k=1$ (equality), and if it holds for $k$ then $\binom{2k+2}{k+1}=\binom{2k}k\frac{2(2k+1)}{k+1}\ge\frac{4^k}{2\sqrt k}\cdot\frac{2(2k+1)}{k+1}\ge\frac{4^{k+1}}{2\sqrt{k+1}}$, the last step being $(2k+1)^2\ge4k(k+1)$. Hence each factor is $\ge\frac{16^k(2k+1)}{4k}>\frac{16^k}2$, and
\[
\begin{gathered}
\frac{c_{2n}}{c_n^4}>2^{-(n-1)}\prod_{k=1}^{n-1}16^k=2^{-(n-1)}\,4^{n(n-1)},\\
\det H_n<2^{\,n-1}4^{-n(n-1)}=2^{-(n-1)(2n-1)}\le4^{-(n-1)^2}.
\end{gathered}
\]
This also gives an elementary conditioning warning. Write $\lambda_{\min},\lambda_{\max}$ for its positive extreme eigenvalues. Since $e_1^{\mathsf T}H_ne_1=1$, the Rayleigh principle gives $\lambda_{\max}\ge1$. Also $\lambda_{\min}^n\le\prod_i\lambda_i=\det H_n$, hence
\[
\frac{\lambda_{\max}}{\lambda_{\min}}\ge(\det H_n)^{-1/n}>4^{(n-1)^2/n}.
\]
Thus the ratio of greatest to least stretch grows without bound. Inverting the least-stretched directions amplifies rounding errors; the determinant estimate and the lower bound on the greatest stretch establish this together.
\end{solution}

\begin{solution}{ex:bridge-translation}
\idea A familiar operation on polynomials becomes transparent in a basis chosen for that operation. Differences shift the binomial basis down one place; translation is the identity plus this shift. The example separates having eigenvalues from having enough eigenvectors.
\writeup (a) Proposition~\ref{prop:deltabinom} gives $\Delta b_0=0$ and $\Delta b_k=b_{k-1}$ for $k\ge1$. Thus the matrix $N$ of $\Delta$ has $1$ immediately above the diagonal and zero elsewhere. The matrix of $S$ is $I+N$, upper triangular with diagonal entries $1$, so $\chi_S=(X-1)^{d+1}$ and $\det S=1$.

(b) Repeatedly shifting a basis element gives $\Delta^{d+1}b_k=0$ for every $k$, whereas $\Delta^db_d=b_0=1\ne0$. The eigenspace for $1$ is $\ker\Delta$, exactly the constants: a nonconstant degree-$r$ polynomial has difference of degree $r-1$ with leading coefficient multiplied by $r\ne0$. This eigenspace has dimension $1<d+1$, so it cannot provide a basis and $S$ is not diagonalisable.

(c) Since $I$ and $\Delta$ commute, the binomial theorem gives $(I+\Delta)^m=\sum_{j=0}^m\binom mj\Delta^j$; terms with $j>d$ vanish, and terms with $j>m$ have zero coefficient. Applying to $P$ gives $P(X+m)=\sum_{j=0}^d\binom mj\Delta^jP(X)$. This is an exact finite translation formula, with no convergence question.
\end{solution}

\subsection*{Notes on the challenges}

(1) The theorem is Perron's, from 1907, extended by Frobenius in 1912 to non-negative matrices that are \emph{irreducible} (the graph of non-zero entries is strongly connected), where the eigenvector stays positive but other eigenvalues of modulus $r$ may appear (a cyclic permutation matrix has all $n$th roots of unity). A proof for positive matrices in the spirit of Exercise~\ref{ex:stochastic}: show that $r=\max\{t\ge0:\ Ax\ge tx\text{ for some }x\ge0,\ x\ne0\}$ is attained, that a maximiser is a positive eigenvector, and that any other eigenvalue has smaller modulus by applying the inequality $|Az|\le A|z|$ entrywise. Horn and Johnson, Chapter~8, is the reference. The damping factor $0.15$ of PageRank makes the Google matrix positive, so that none of the subtleties of the reducible case arise.

(2) Normalise a Hadamard matrix so that its first row is all ones (multiply columns by $\pm1$); orthogonality of rows $1$ and $2$ says row $2$ has $n/2$ entries $+1$, so $n$ is even; then counting the columns where rows $2$ and $3$ agree in sign, using orthogonality of the three pairs, gives $n\equiv0\pmod4$ for $n\ge3$. The conjecture that Hadamard matrices exist for every $n\equiv0\pmod4$ is from Paley's paper of 1933 and remains open; order $428$ was constructed in 2004 (the paper appeared in 2005). In August 2026, Alp\"oge, Voinov and Reynolds-Haertle announced the twelve previously unknown orders below $2000$, including $668$; Egan and coauthors verified them in their October 2026 survey, which lists $2396$ as the smallest still-open order.

(3) For $n=2$, writing the nilpotent $N=\begin{pmatrix}0&1\\0&0\end{pmatrix}$ as $S_1S_2$ is already a nice puzzle (try $S_1=\begin{pmatrix}0&1\\1&0\end{pmatrix}$, and $S_2$ follows); the general theorem is Frobenius's, from 1910, and the clean proof uses the rational canonical form of Volume~3: a companion matrix is visibly a product of two symmetric matrices, and the property passes to direct sums and to similar matrices (if $A=S_1S_2$ then $P^{-1}AP=(P^{-1}S_1P^{-\mathsf T})\cdot(P^{\mathsf T}S_2P)$, both symmetric). A.~J.~Bosch's note in the \emph{Monthly} gives a two-page account.

\begin{solution}{pb:tanroots}
\idea Near $a=(k+\frac12)\pi$ the tangent is huge, so $\tan x=x$ forces $x$ just below $a$; writing $x=a-\varepsilon$ turns the equation into $\tan\varepsilon=\frac1{a-\varepsilon}$, which can be iterated.

\writeup (a) On $I_k=\left]k\pi-\frac\pi2,k\pi+\frac\pi2\right[$ the function $h(x)=\tan x-x$ has $h'=\tan^2x\ge0$, vanishing only at $k\pi$, so $h$ is increasing, from $-\infty$ to $+\infty$: exactly one root.
Since $h(k\pi)=-k\pi<0$, the root is in $\left]k\pi,k\pi+\frac\pi2\right[$.

(b) Put $x_k=a-\varepsilon$ with $0<\varepsilon<\frac\pi2$. Then $\tan(a-\varepsilon)=\cot\varepsilon$, so the equation is $\tan\varepsilon=\frac1{a-\varepsilon}$. Since $a-\varepsilon>a-\frac\pi2\ge\frac a3$, $\varepsilon<\tan\varepsilon\le\frac3a$. Then
$\frac1{a-\varepsilon}=\frac1a\bigl(1+\frac\varepsilon a+O(\frac{\varepsilon^2}{a^2})\bigr)=\frac1a+\frac1{a^3}+O(a^{-5})$, using $\varepsilon=\frac1a+O(a^{-3})$, itself obtained from $\varepsilon=\arctan\frac1{a-\varepsilon}=\frac1a+O(a^{-3})$. Finally
$\varepsilon=\arctan y=y-\frac{y^3}3+O(y^5)$ with $y=\frac1a+\frac1{a^3}+O(a^{-5})$ gives $\varepsilon=\frac1a+\frac2{3a^3}+O(a^{-5})$.

(c) The expansion is in powers of $\frac1a$, and $a=\frac{3\pi}2=4.71$ already for $k=1$; the next term, computed numerically, is $-\frac{13}{15a^5}$, about $-0.0004$, which is the
observed error.

\tools{Chapter~\ref{ch:calculus}: the derivative, and the Taylor expansion of $\arctan$.}
\leads These roots are the zeros of $\frac{\sin x-x\cos x}{x^3}=\frac13-\frac{x^2}{30}+\cdots$. A product expansion like Euler's sine product then gives $\sum_k\frac1{x_k^2}=\frac{1/30}{1/3}=\frac1{10}$, which the
numerical evaluation confirms. In practice such expansions are excellent starting points for Newton's method.
\end{solution}

\section*{Sources and further reading}

\begin{itemize}\raggedright
\item R.~Egan, I.~S.~Kotsireas, S.~London and P.~\'O~Cath\'ain, \emph{Recent progress in the construction of Hadamard matrices}, arXiv:2610.27175, 19 October 2026. Verification of the new orders, including $668$, and the remaining frontier.
\item S.~Axler, \emph{Linear Algebra Done Right}, 4th ed., Springer, 2024 (open access). Eigenvalues, inner products and the spectral theorem, with determinants at the end.
\item P.~D.~Lax, \emph{Linear Algebra and Its Applications}, 2nd ed., Wiley, 2007. Short, deep, and full of applications: Rayleigh quotients, positive definite and positive matrices.
\item R.~A.~Horn and C.~R.~Johnson, \emph{Matrix Analysis}, 2nd ed., Cambridge University Press, 2013. The reference for everything in this chapter and much more (Chapter~4 for Rayleigh quotients, Chapter~8 for Perron--Frobenius).
\item T.~Muir, \emph{The Theory of Determinants in the Historical Order of Development}, vol.~1, Macmillan, 1906. Seki's manuscript of 1683, Leibniz's letter to l'H\^opital of 1693, Cramer's rule (1750) and Cauchy's memoir of 1812 (printed in \emph{J.\ \'Ecole Polytechnique} 10, 1815), all quoted at length.
\item A.~Cayley, \emph{A memoir on the theory of matrices}, Philos.\ Trans.\ R.\ Soc.\ London 148 (1858), 17--37. Matrices as objects, and the theorem of Theorem~\ref{thm:cayleyhamilton}, verified for $n\le3$; the general proof is G.~Frobenius, \emph{Ueber lineare Substitutionen und bilineare Formen}, J.\ reine angew.\ Math.\ 84 (1878), 1--63. The word ``matrix'' is from J.~J.~Sylvester, Philos.\ Mag.\ 37 (1850), 363--370.
\item A.-L.~Cauchy, \emph{Sur l'\'equation \`a l'aide de laquelle on d\'etermine les in\'egalit\'es s\'eculaires des mouvements des plan\`etes}, Exercices de math\'ematiques 4 (1829). Real eigenvalues of a symmetric form, before matrices. The determinant of Exercise~\ref{ex:cauchydet} is in his \emph{M\'emoire sur les fonctions altern\'ees et sur les sommes altern\'ees}, Exercices d'analyse et de physique math\'ematique 2 (1841), 151--159.
\item D.~Hilbert, \emph{Ein Beitrag zur Theorie des Legendre'schen Polynoms}, Acta Math.\ 18 (1894), 155--159. The matrix $H_n$; and M.-D.~Choi, \emph{Tricks or treats with the Hilbert matrix}, Amer.\ Math.\ Monthly 90 (1983), 301--312, for its exact inverse and why it is so badly conditioned.
\item J.~Hadamard, \emph{R\'esolution d'une question relative aux d\'eterminants}, Bull.\ Sci.\ Math.\ 17 (1893), 240--246. The inequality of Exercise~\ref{ex:hadamard}; R.~E.~A.~C.~Paley, \emph{On orthogonal matrices}, J.\ Math.\ Phys.\ 12 (1933), 311--320, for the conjecture; H.~Kharaghani and B.~Tayfeh-Rezaie, \emph{A Hadamard matrix of order 428}, J.\ Combin.\ Des.\ 13 (2005), 435--440.
\item O.~Perron, \emph{Zur Theorie der Matrices}, Math.\ Ann.\ 64 (1907), 248--263; G.~Frobenius, \emph{\"Uber Matrizen aus nicht negativen Elementen}, Sitzungsber.\ Preuss.\ Akad.\ Wiss.\ (1912), 456--477. Challenge (1).
\item A.~J.~Bosch, \emph{The factorization of a square matrix into two symmetric matrices}, Amer.\ Math.\ Monthly 93 (1986), 462--464. Challenge (3), with Frobenius's 1910 theorem.
\item S.~Brin and L.~Page, \emph{The anatomy of a large-scale hypertextual Web search engine}, Computer Networks and ISDN Systems 30 (1998), 107--117; and K.~Bryan and T.~Leise, \emph{The \$25{,}000{,}000{,}000 eigenvector: the linear algebra behind Google}, SIAM Review 48 (2006), 569--581. PageRank, as in Exercise~\ref{ex:stochastic}.
\item C.~F.~Gauss, \emph{Theoria motus corporum coelestium}, Perthes and Besser, 1809. Least squares and the orbit of Ceres; Legendre had printed the method first, in 1805, and the quarrel over priority is told in S.~M.~Stigler, \emph{The History of Statistics}, Harvard University Press, 1986, Chapter~1.
\item R.~Simson, \emph{An explication of an obscure passage in Albert Girard's Commentary upon Simon Stevin's works}, Philos.\ Trans.\ 48 (1753), 368--377. The identity $F_{n+1}F_{n-1}-F_n^2=(-1)^n$, which Cassini had stated in the \emph{Histoire de l'Acad\'emie royale des sciences} for 1680; and W.~Weaver, \emph{Lewis Carroll and a geometrical paradox}, Amer.\ Math.\ Monthly 45 (1938), 234--236, for the $64=65$ puzzle in Carroll's papers.
\end{itemize}
